% Автоматически перенесено из ГЛАВА_3_ВЕРСИЯ_3.docx

\chapter{Система доказательств в Каталонии XII--XIII~вв.}

Судебное доказательство в Каталонии XII--XIII~вв. представляет собой одну из наиболее показательных областей для изучения правовой культуры эпохи, поскольку именно в способах установления истины особенно ясно проявляются представления общества о праве, справедливости, авторитете и допустимых формах разрешения конфликта. Средневековый суд не располагал единым и универсальным набором доказательств в современном смысле слова: истина устанавливалась через сочетание сакральных практик, устных заявлений сторон, свидетельских показаний, письменных актов, репутации участников процесса и, позднее, элементов розыскного расследования. Поэтому анализ средств доказательства позволяет увидеть не только технику судопроизводства, но и более глубокий уровень --- то, каким образом право связывало факт, слово, социальный статус и коллективное признание.

Для Каталонии рассматриваемого периода особенно важно, что речь идет о правовом пространстве, в котором сосуществовали различные нормативные традиции. «Барселонские обычаи», городские привилегии, «Обычаи Льейды», «Обычаи Тортосы», а также зависимые от них локальные сборники отражают не только локальное разнообразие, но и постепенное изменение самой логики судебного доказательства. В одних текстах ещё преобладают более ранние формы, тесно связанных с ордалией, судебным поединком и статусной природой суда; в других --- отчётливо усиливается стремление к письменной фиксации и процессуальной определённости.

В этом смысле эволюция форм судебного доказательства позволяет проследить, в каком направлении развивалась правовая культура каталонского средневековья в области форм удостоверения истины. Именно поэтому в центре внимания настоящей главы находится не столько формальная классификация или типология судебного доказательства для обозначенных периода и региона, сколько вопрос, какие социальные, религиозные и процессуальные предпосылки делали то или иное средство предпочтительным в каждом конкретном случае.

Особое значение для такого анализа имеет антропологический подход к правовым текстам. В его перспективе доказательство понимается не просто как техническая процедура, а как форма социального действия, укоренённая в представлениях общины о дозволенном и выходящем за его рамки. Средневековый суд имел дело с фактами и документами, уже включёнными в систему социальных отношений и коллективных ожиданий. Поэтому испытание, клятва, свидетельство или добрая слава значимы не только как отдельные институты, но и как разные способы превращения социально признанного знания в юридически обязательную истину. В этом отношении судебное доказательство оказывается одной из важнейших сфер, в которых право выступает механизмом социализации индивида и одновременно средством воспроизводства самой общины.

Центральная задача этой главы настоящего исследования --- показать, каким образом в каталонских правовых памятниках XII--XIII~вв. были организованы основные способы доказательства. В центре внимания находятся, с одной стороны, сакральные и ритуальные виды доказательства --- прежде всего ордалии, судебный поединок и клятва, --- а с другой стороны, формы, связанные с возрастающей процессуальной дисциплиной: репутационные критерии, свидетельские показания, допрос, презумпции и документы. Такой порядок рассмотрения позволяет проследить внутреннюю логику их соотношения и изменения.

Первая группа рассматриваемых форм связана с тем, что истина в раннем и высоком Средневековье мыслилась как открываемая через сакральный знак, ритуально оформленное действие или публично подтверждённое слово. Вторая группа форм --- репутация, свидетельство, допрос, документ--- позволяет увидеть, как в судебный процесс всё активнее включались механизмы социальной проверки и процедурного контроля.

Таким образом, в настоящей главе судебное доказательство рассматривается как пространство, в котором пересекаются сакральное, социальное и процессуальное измерения права. Через анализ ордалий, клятвы, репутации, свидетельских показаний и розыскного расследования представляется возможным показать, что юридическая истина в Каталонии XII--XIII~вв. производилась не одним каким-либо универсальным способом, а в результате взаимодействия различных правовых логик. Именно в этой сфере особенно ясно видно, как право переходило от более ранних, персонализированных и статусных форм к более институционально организованным способам разрешения конфликта, не утрачивая при этом своей глубокой связи с коллективной памятью и социальными представлениями общины.

\section{Сакральное испытание и судебный поединок}

В историографии ордалии обычно определяются как особый вид судебного доказательства, при котором виновность или невиновность устанавливалась посредством физического испытания, понимавшегося как проявление Божьей воли; отсюда и распространенное латинское обозначение \emph{«Божий суд»} (лат. \emph{--- }\emph{judicium}\emph{ Dei}). В самом общем виде исследователи различают двусторонние и односторонние ордалии. К двусторонним относился прежде всего судебный поединок, при котором в испытании участвовали обе стороны спора или их представители, а исход борьбы воспринимался как указание на правоту одной из сторон\footnote{\emph{Тогоева}\emph{ О.И. }Ордалия // Православная энциклопедия. 2023. Т.~53. С.~104--106.}.

Для рассматриваемого каталонского материала наибольшее значение имеют, с одной стороны, судебный поединок (batalia) как двусторонняя ордалия, а с другой --- односторонние испытания раскаленным железом, кипятком и холодной водой. Испытания железом и кипятком были построены по сходной логике: обвиняемый подвергался физическому воздействию, а затем через несколько дней оценивалось состояние тела; отсутствие явных следов повреждения толковалось в его пользу. Ордалия холодной водой, напротив, предполагала немедленно наблюдаемый результат: связанного обвиняемого опускали в воду, и по его поведению в воде судили об исходе испытания. Такая типология показывает, что различия между видами ордалий определялись не только используемым элементом, но и самой процессуальной конструкцией испытания: в одном случае правота стороны выявлялась через состязание двух участников, в другом --- через телесную проверку одного обвиняемого.

В каталонских источниках судебное испытание, обозначаемое как «Божий суд» (лат. --- \emph{judicium}\emph{ Dei}), предстаёт одним из установленных способов судебного доказательства. Его применение определялось характером разбираемого дела, процессуальным положением сторон и их социальным статусом и соотносилось с другими средствами подтверждения правоты --- клятвой, свидетельскими показаниями и письменными актами. Связь ордалии с клятвой не сводилась к альтернативе между двумя самостоятельными процедурами. Клятва могла входить в подготовку судебного испытания, сопровождать его или предшествовать ему, тогда как обращение к «Божьему суду» позволяло передать разрешение спора сакральной инстанции в тех случаях, когда предусмотренная процедура требовала испытания, способного придать судебному решению дополнительную убедительность и обязательность.

Современные исследования судебной практики Каталонии XII--XIII~вв. показывают, что ордалии являлись полноценным инструментом правосудия: по установленному церковному обряду проверка помогала уточнить или подтвердить обстоятельства дела\footnote{\emph{Baldwin J. }The Intellectual Preparation for the Canon of 1215 against Ordeals // Speculum. Cambridge, 1961. Vol.~36. P.~613--636;\emph{ Barthélemy D. }Présence de l’aveu dans le déroulement des ordalies (IXe-XIIIe siècle) // Publications de l’École française de Rome. École Française de Rome, 1986. Т.~88, №~1. P.~191--214; \emph{Idem. }Diversité des ordalies médiévales // Revue historique. Paris, 1988. Vol.~280. №~1. P.~3--25; \emph{Bartlett R. }Trial by Fire and Water: The Medieval Judicial Ordeal. Oxford, 1986; 2 vol.; \emph{Carbasse J.-M., Vielfaure P. }Histoire du droit pénal et de la justice criminelle. Presses Universitaires de France, 2014. 544 p. Довольно подробная историографическая панорама, посвящённая судебному поединку на Пиренейском полуострове (в частности, в Леоно-Кастильской традиции) представлена в работе А. Марея: \emph{Марей А.В.} Божий суд и человеческая справедливость: судебный поединок в Леоне и Кастилии XI-XIII~вв. Москва: Союзник, 2011. 54 с.}.

\subsection{Виды ордалий и сфера их применения}

В «Барселонских обычаях» ордалия в форме испытания водой упоминается уже в 1-й статье в одном ряду с очистительной клятвой и судебным поединком: \emph{«…клятвой, или поединком, или холодной либо горячей водой…»}\footnote{Usatges de Barcelona... Us. 1 (1--2): «…per sacramentum, vel per batalliam, vel per aquam frigidam sive calidam…». P.~48--49.}. При этом проведение любого из этих действий предполагало произнесение клятвенной формулы: \emph{«Все злодеяния, которые я учинил против тебя, я совершил по своему прямому праву и по причине твоей небрежности, поэтому не должен их перед тобою исправлять. Клянусь Богом и его святыми»}\footnote{Ibid.: «Iuro ego tibi quod hec malefacta, que tibi habeo facta, sic ea tibi feci ad meum directum et in tuo neglecto, quod ego tibi illa emendare non debeo, per Deum et per hec sancta». P.~48--49.}. Таким образом, роль клятвы среди других сакральных способов доказательства выделяется на фоне ордалии и судебного поединка. И лишь когда клятва оказывалась недостаточной для установления «истины», согласно этой норме «Барселонских обычаев», могли применяться испытания водой и судебные поединки. Это свидетельствует о том, что в XI--XII~вв., когда были составлены «Барселонские обычаи», ордалия уже постепенно вытеснялась другими формами сакрального доказательства.

Это обстоятельство подтверждается и указанием самих «Барселонских обычаев» на то, что в данном случае они фиксировали лишь сложившуюся правовую традицию, существовавшую в предшествующий период: \emph{«До того, как }\emph{«Обычаи» были дан}\emph{ы…»}\footnote{Ibid.: «ANTEQVAM VSATICI fuissent missi…». P.~48--49.}, от неё пытались отчасти дистанцироваться составители этого кодекса. Показательно, что 90-я статья «Барселонских обычаев» разрешает проведение ордалии только в случае \emph{«если истину невозможно установить иным способом»}\footnote{Usatges de Barcelona... Us. 90 (113): « VERE INDEX aliter non erit nisi hoc quod indicauerit ad uerum traxerit …». P.~126--127.}. Стоит также отметить, что вслед за «Вестготской правдой»\footnote{LV VI.1.3. Каким образом судья должен расследовать дело посредством испытания кипятком: «Нам стало известно о~многих жалобах, и~мы узнали, что проводятся пытки по~делам [стоимостью до] трехсот солидов и~претерпевается много зла от~свободных людей. И~потому, для~наилучшего порядка, мы устанавливаем и~приказываем, чтобы, хотя~бы дело по~какому-либо обвинению и~было незначительным, обвиняемый все же должен быть принужден судьей к~испытанию кипятком; и~когда обстоятельства дела прояснятся, пусть судья не~усомнится подвергнуть его пытке, и~после того как тот даст свое признание, пусть он будет подвергнут наказанию, определенному предыдущим законом. Если~же он выйдет невиновным из~испытания кипятком, то пусть истец не~несет никакого наказания. Тот~же порядок должен соблюдаться и~относительно подозрительных людей, пришедших свидетельствовать» C.~335--336, 646--647.} «Барселонские обычаи» содержат упоминания исключительно об испытаниях водой (упоминания об ордалии калёным железом, пищей, чудом и др. в них отсутствуют), однако в них, в отличие от «Книги приговоров», процедура проведения ордалии кипятком или холодной водой в них не описывалась. Вероятно, предполагалось, что она была общеизвестной и не требовала процессуального разъяснения; указывалось лишь, что крестьянки, подозреваемые в супружеской неверности, \emph{«...должны сами, своими руками, проходить испытание кипятком»}\footnote{Usatges de Barcelona... Us. 89 (112): «...uxores rusticorum manibus propriis per caldariam». P.~126--127.}.

Как можно понять из 1-й статьи «Барселонских обычаев», испытание водой применялось прежде всего в случае уголовных преступлений (убийства, нарушения присяги, супружеская неверность и т. д.). В то же время 50-я статья «Барселонских обычаев» предполагала использование ордалии кипятком для крестьян-баккалариев в имущественных спорах: \emph{«Прочим же крестьянам, называемым }\emph{баккалариями}\emph{, доверяют на основании клятвы лишь до суммы в четыре }\emph{манкуса}\emph{ валенсийского золота; всё, что свыше, в отношении чего они приносят клятву, должно быть доказано посредством испытания кипятком»}\footnote{Ibid. Us. 50 (53): «DE ALIIS NAMQVE rusticis qui dicunty challarii, credantur per sacramentum usque ad. ba. mancusos auri Valencie. Deinde quicquid iurent pe examen caldarie demonstrent». P.~88--89.}\emph{. }Таким образом, «Божий суд» оказывался встроен в социальную иерархию, организованную по принципу выбора наиболее адекватного происхождению подсудимого или свидетеля способа доказательства. Испытания кипятком не могли применяться в отношении обвиняемых знатного или рыцарского происхождения.

\subsection{Социальный статус участников судебного испытания}

Нормы «Барселонских обычаев», регулирующие судебный поединок, в ещё большей степени обнаруживают его зависимость от социальной стратификации. Вместе с тем в них последовательно проводится принцип равенства сторон, без которого сама процедура теряла бы легитимность. Рассмотрим статью «Барселонских обычаев»: \emph{«…или очиститься от обвинения в }\emph{баузии}\emph{ клятвой и поединком с равным себе, который по происхождению и чести соответствует его достоинству...»}\footnote{Ibid. Us. 42 (45): «... aut expiare de de bauzia per sacramentum et per bataiam ad suum parem, qui de genere et de honore sit de suo valore…». P.~84--85.}. Таким образом, требование поединка с «равным себе» означало, что юридическая сила такого доказательства была неотделима от социально признанного равенства участников. В данном случае речь идёт об обвинении в тяжком преступлении, баузии\footnote{\textbf{Баузия} (лат. bausia) --- вероломство, коварное нарушение верности или мира, то есть тяжкое деяние, совершённое против лица, которому была обязана особая верность или безопасность. В каталонских юридических текстах это понятие обычно сближается с изменой и неверностью сеньору, но обычно передаёт не столько государственную измену в узком смысле, сколько вероломное посягательство на лицо, честь или безопасность при нарушении должной верности. (см. \emph{Zambrana Moral P.} Revenja privada i revenja de la sang en el dret penal espanyol medieval: estat de la qüestió en els inicis del 2005 // Revista de Dret Històric Català. 2006. Т.~5. P.~128; \emph{Alvarado Planas J. }Algunas observaciones sobre la influencia germánica en el vocabulario jurídico-institucional de la España medieval // Glossae: European Journal of Legal History. Instituto de Estudios Sociales, Políticos y Jurídicos, 2015. №~12. P.~46).}; избавление от такого обвинения требовало не только проведения судебного поединка, но и принесения очистительной клятвы. Таким образом, согласно «Барселонским обычаям», при проведении т.н. «очистительного» поединка и обвинении в тяжких преступлениях был важен не только исход судебного поединка (победа или проигрыш), но и социальный статус его участников.

В то же время далее в тексте 42 статьи «Барселонских обычаев» имеются сведения о возможной выплате участникам в случае победы в судебном поединке: \emph{«…чтобы он получил столько же в случае победы, сколько потерял бы в случае поражения»}\footnote{Usatges de Barcelona… Us. 42 (45): «…ut tantum accipiat si vincerit quantum perderet si vinceretur». P.~84.}. Эта формула показывает, что судебный поединок мыслился не только как вид доказательства, но и как процедура с заранее определёнными имущественными последствиями. Победа и поражение здесь построены симметрично: исход поединка влечёт для стороны либо выигрыш, либо проигрыш в одном и том же объёме (лат. --- damnum et produm). Отдельно в 23-й статье «Барселонских обычаев» оговаривались размеры залогов для пеших и конных поединков, из которых предполагалась выплата компенсаций за возможные увечья, потерю коня и оружия: \emph{«Если судебный поединок назначен, то прежде, чем он будет подтверждён клятвой, он должен быть обеспечен залогами: если его надлежит вести рыцарям --- в размере двухсот унций валенсийского золота, а если пешим воинам --- в размере ста; для того, чтобы тому, кто победит, был возмещён вред, который он понесёт в поединке --- как телесный, так и в коне либо в оружии, --- и чтобы он получил то, из-за чего поединок был начат, а также все издержки, которые понесёт вследствие этого поединка; и должно быть определено то, что он получит с того, кто будет побеждён»}\footnote{Ibid. Us. 23 (27): «BATAIA IVDICATA antequam sit iurata, si per milites debet esse facta, per .CC. uncias auri Valen. cie sit per pignora firmata; et si per pedones, sit firmata per .c., propter hoc ut ad illum qui uicerit sit emendatum malum quod in bello acceperit, tam in corpore quam in cauallo siue in armis, et assequatur hoc pro quo bellum factum erit et omnes missio. nes quas per illud bellum fecerit; et diffinitum illud quod acceperit ille qui uictus fuerit». P.~70.}.

Сумма в 200 унций валенсийского золота была по меркам самих «Барселонских обычаев» исключительно крупной. Сопоставление с другими денежными величинами, упоминаемыми в тексте\emph{ }«Барселонских обычаев» --- 4 манкуса, 5, 20 и 100 унций валенсийского золота, --- показывает, что требуемый для рыцарского поединка залог относится к максимальным суммам и потому должен рассматриваться как форма чрезвычайного процессуального обеспечения\footnote{\emph{ Balaguer A.M. }Història de la moneda dels comtats catalans. Barcelona: Institut d’Estudis Catalans, 1999. P.~55.}. Тем самым поединок выступал как строго оформленный правовой механизм, соединявший доказательство, санкцию и компенсацию.

В «Барселонских обычаях» женские ордалии были описаны в 89-й статье как система испытаний, напрямую зависевшая от социального положения семьи. Эта статья связана с процессуальными особенностями обвинения в супружеской неверности: \emph{«Мужья вправе обвинять своих жён в прелюбодеянии даже по одному подозрению, а жёны должны оправдываться от такого обвинения собственным заверением, клятвой и поединком}»\footnote{Usatges de Barcelona… Us. 23 (27): «MARITI VXORES suas reptare possint de adulterio, eciam per suspicionem, et ille debent se expiare inde per illarum aueramentum, per sacramentum et per bataliam». P.~124.}.

Для жён рыцарей предусматривалась форма, связанная с судебным поединком через представителей:  \emph{«жёны рыцарей [очищаются.---А.К.] клятвой и, сверх того, через рыцаря»}\footnote{Ibid. Us. 89 (112): «uxores militum per sacramentum et desuper per militem». P.~126.}. Для женщин более низкого социального статуса предполагалось проведение судебного поединка через пешего представителя: \emph{«жёны горожан и знатных }\emph{байюлов}\emph{ [очищаются.---А.К.] через пешего бойца»}\footnote{Ibid. Us. 89 (112): «uxores civium et burgencium et nobelium bajulorum per pedonem». P.~126.}. Напротив, для крестьянок устанавливалась личная ордалия кипятком: \emph{«жёны крестьян должны сами, своими руками, проходить испытание кипятком} \footnote{Ibid. Us. 89 (112): «uxores rusticorum manibus propriis per calderiam». P.~126.}. Тем самым источник ясно показывал, что участие женщин в судебном испытании регулировалось в соответствии с их положением в социальной иерархии: жёны рыцарей, горожан и знатных байюлов участвовали в судебном поединке через представителей, тогда как крестьянки должны были лично пройти испытание кипятком. Исход процедуры также формулировался предельно определённо: \emph{«Если жена победит, муж должен удержать её при себе с честью…»}\footnote{Ibid. Us. 89 (112): «Si vicerit uxor, retineat eam vir suus honorifice…» P.~126.}; если же она проигрывала, то \emph{«пусть перейдёт во власть мужа со всем, что имеет»}\footnote{Ibid. Us. 89 (112): «Si autem victa fuerit, veniat in manu mariti sui cum cunctis que habuerit». P.~126.}. Женские ордалии в «Барселонских обычаях» выступали частью системы доказательств, в которой выбор между поединком через представителя и личным испытанием кипятком зависел от социального положения женщины и её семьи.

«Барселонские обычаи» также содержали нормы, допускавшие замену судебного поединка клятвой при определённых обстоятельствах: «\emph{Престарелому рыцарю, не способному защищаться собственными силами, а также бедняку, который не в состоянии снарядиться для поединка, следует доверять по его клятве в пределах пяти унций валенсийского золота»}\footnote{Ibid. Us. 51 (54): «SENEX MILES qui non potuerit se defendere per semedipsum, uel pauper qui non potest se preparare de bello, credatur per sacramentum usque ad ‚V. uncias auri Valencie». P.~88.}. Тем самым «Барселонские обычаи» признавали, что доступ к поединку зависел не только от статуса стороны, но и от её физических и материальных возможностей. Вместе с тем речь не шла о полном освобождении от поединка как способа доказательства: клятва признавалась достаточной лишь в пределах строго установленной суммы, тогда как для прочих рыцарей обычай, напротив, сохранял обязанность защищать свои показания в поединке собственными руками против равного себе. Поэтому норма показывает не отказ от судебного поединка, а его гибкую адаптацию к возможным конкретным обстоятельствам дела.

В целом «Барселонские обычаи» свидетельствуют о том, что судебный поединок и ордалии в Каталонии XII~в. представляли собой довольно подробно регламентированные виды доказательства. Это подтверждается и материалами грамот, в которых эти способы доказательства выглядят весьма частотными. Их применение подчинялось сословной логике, зависело от статуса и процессуальной способности сторон и потому раскрывает глубокую связь судебного доказательства с социальной структурой каталонского общества.

\subsection{Ордалия и судебный поединок в актовой практике XI--XII~вв.}

Количественный анализ документов, опубликованных в 1-м и 2-м томах сборника «Правосудие и разрешение конфликтов в средневековой Каталонии»\footnote{Justícia i resolució... 2018.; Justícia i resolució... 2024.}, показывает, что и в XI, и в XII~вв. среди ордалий и смежных доказательных процедур преобладал судебный поединок, тогда как односторонние ордалии встречались заметно реже. В выборке из 552 документов IX--XI~вв. судебный поединок фиксируется исключительно в документах XI~в. как минимум в 10 из 15 документов с упоминанием об ордалиях, тогда как испытание кипятком --- в 2 случаях, холодной водой --- в 1, а испытание калёным железом не засвидетельствовано вовсе, что соответствует ситуации, описанной в «Барселонских обычаях». В минимально выявляемом по тексту тома II наборе из 19 документов с упоминанием ордалии в XII~в. поединок встречается не менее чем в 6 документах, ордалия кипятком --- в 4, а холодной водой и калёным железом --- по 1 разу соответственно. В обеих выборках присутствует по одному «смешанному случаю». Это позволяет предположить, что на всём протяжении XI--XII~вв. именно судебный поединок оставался наиболее значимой формой «Божьего суда», тогда как прочие разновидности играли весьма ограниченную роль в каталонской правовой культуре.

«Божий суд» и в 1-й пол. XIII~в. оставался одним из средств установления виновности/невиновности в ходе судебных процессов, однако носил необязательный характер. Так, в привилегии 1211~г., пожалованной г.~Балагеру Пере I (1196--1213), в качестве «вольности» встречается «освобождение» горожан от ордалий калёным железом, холодной и горячей водой: \emph{«…и да будут они всегда свободны, неприкосновенны и полностью освобождены от суда или очистительного испытания калённым железом и горячей либо холодной водой, так что в такого рода судах или очистительных процедурах должна иметь силу лишь их простая клятва, за }\emph{исключением, однако, вышеупомянутых очистительных процедур между соседями»}\footnote{Concessió de privilegi atorgat per Pere I a la ciutat de Balaguer de poder celebrar mercat el dissabte i una fira de vuit dies a mitjans d’agost. Lleida, 16 de setembre de 1211. Arxiu Comarcal de la Noguera (=ACN). Fons Ajuntament de Balaguer. UDS: 2001 // \emph{Pedrol M. }\emph{e}\emph{t}\emph{~}\emph{al. }El mercat de Balaguer a través dels documents de l’Arxiu Comarcal de la Noguera. (Quadern didàctic ACN, 2008)~(транскрипция документа---А.К.): «… semper sint liberi et immunes et penitus exempti a iudicio sive purgatione candentis ferri et aque calide et frigide, ita scilicet ut in huiusmodi iudiciis sive purgationibus simplex solummodo iuramentum ipsorum teneat, salvis tamen vicinalibus purgationibus supradictis». P.~[3].}.

\subsection{Ограничение судебных испытаний в городском праве XIII~в.}

После 1215~г., когда IV Латеранский собор запретил духовным лицам участвовать в проведении ордалий, их применение оказалось существенно затруднено\footnote{XVIII. De iudicie sanguinis et duelli cliricis interdicto // \emph{Mansi J.D.}~Sacrorum conciliorum nova et amplissima collectio... Venetiis: apud Anitonio Zatta. 1778. T.~22. Col. 1006--1007. Col. 1006--1007.}. Однако из каталонской правовой культуры «Божий суд» не исчез одномоментно, хотя сфера его применения постепенно сужалась.

Например, «Обычаи Льейды», составленные легистом Гийомом Ботетом для этого города в 1227--1228~гг., запрещали проведение любых ордалий в принципе: \emph{«Кроме того, устанавливается, что мы не должны подвергаться вместе с ними ни судебному поединку, ни испытанию железом, ни какому-либо другому виду суда, ни испытанию водой}» \footnote{Costumbres de Llerida. [27] // \emph{Mar C.J.}~Bujaraloz~: VIII centenario de su fundacion y epoca de su pertenencia a la Orden de San Jorge de Alfama 1205-1230 / \emph{C.J.~Mar}. Caspe: Centro de Estudios Comarcales del Bajo Aragon-Caspe de la Institucion “Fernando el Catolico”, 2006: «Item addunt quod bataillam non teneamur facere cum eis per hominem vel ferrum nec per aliud iudicium, nec per aquam». P.~224.}. Данная статья закрепила освобождение от судебного поединка, испытания железом, водой и иных сходных форм «Божьего суда». В этом отношении она показательна для эволюции каталонской городской правовой культуры XIII~в.: ордалии и поединок были изъяты из числа обычных и обязательных процедур доказательства. Норма, таким образом, свидетельствует о сужении сферы применения архаических форм доказательства и о постепенном усилении иных, более процессуально контролируемых способов установления обстоятельств дела.

В то же время, например, 9 статья «Обычаев Тарреги» (1242), пожалованных этому городу Жауме I Завоевателем (1213--1276), закрепляла право проведения испытания кипятком, холодной водой, калёным железом или судебного поединка исключительно по взаимному соглашению сторон: \emph{«Никого не }\emph{дóлжн}\emph{o}\emph{ принуждать ни к поединку, ни к испытанию горячей, либо холодной водой, либо раскалённым железом, если на то не будет воли сторон»}\footnote{Consuetudines ville Tarege… [9]: «Stare non teneatur aliquis ad batalliarri neque tornas aque calide vel frigide vel ferri calidi nisi de voluntate partium fuerit». P.~438.}. Издатель этого памятника Дж.М.~Фонт-Руис отмечал также, что практика подобного ограничения ордалий была распространена ещё до выхода XVIII канона IV Латеранского собора\footnote{Consuetudines ville Tarege… P.~439.}. При этом составители «Кутюмов Перпиньяна» умалчивали даже о возможности применения ордалии на территории города, ограничившись лишь замечанием о том, что для судебных разбирательств не могли использоваться «Барселонские обычаи» и «Вестготская правда»\footnote{Les Coutumes de Perpignan. [I]~: «Homines Perpinyani debent placitare et iudicari per consuetudines ville, et per iura ubi consuetudines defficiunt, et non per usaticos Barchinone, neque per legem goticam, quia non habent locum in villa Perpinyani». P.~567.}. Всё это подразумевало отказ от практики ордалий в пользу более рациональных способов установления истины.

В 1270-е гг. обе редакции «Обычаев Тортосы» демонстрировали жёсткое ограничение всех видов ордалий (в том числе и судебного поединка) в статье 1.1.12: \emph{«Издавна существовал общий обычай, согласно которому никто --- ни сеньор, ни иной человек --- не может вызвать кого-либо на какой бы то ни было судебный поединок, пеший или конный, ни на испытание раскалённым железом, ни на испытание горячей или холодной водой, поскольку в Тортосе, равно как и в её округе, судебный поединок не допускается»}\footnote{Costums de Tortosa… 1.1.12: «Totz temps jo general custum e és que nul hom, seynor ne aitre, no pot reptar negun que a nula batayla, de peu, ne de cavai, ne de ferre calt, ne a ayga calda o freda la port, per so corn batala noa en Tortosa ne yjo, ne en sos ténnens». P.~11.}\emph{. }43-я статья королевской привилегии «Знатные мужи Барселоны постановили…», пожалованной Пере~III Великим г.~Барселоне в 1283--1284~гг., существенно ограничивала область применения судебного поединка как средства установления виновности или невиновности обвиняемого рядом тяжких преступлений: нарушения перемирия, баузии или предательства\footnote{Recognoverunt proceres [XLIII]…: «…nemo potest aliquem reptare de batallia in curia Barchinone… nisi… de treugis fractis vel bausia vel traditione». P.~589.}. Это обстоятельство вводило судебный поединок в сферу преступлений «государственных», направленных на подрыв перемирия Божьего, клятв верности и т. д.

\section{Судебная клятва: ритуал и процессуальная функция}

\subsection{Место клятвы среди средств доказательства}

Судебная клятва представляет собой правовой феномен, находящийся на пересечении нормативно-правовых, моральных и религиозных регулятивных систем. М.~Фуко называл средневековую судебную клятву \emph{«древним магически-религиозным дознанием»}\footnote{\emph{Фуко М. }Истина и правовые установления // Интеллектуалы и власть: Избранные политические статьи, выступления и интервью. Ч. 2 / пер. с франц. И. Окуневой под общ. ред. \emph{Б.М.~Скуратова}. М., 2005. С.~83.}. В условиях средневекового общества, в котором право ещё не было строго формализовано, а письменные документы часто уступали место устным показаниям и ритуальным действиям, клятва становилась одним из важных инструментов подтверждения слов и намерений.

В интересующий нас период клятвы и присяги активно использовались не только в суде, но и в повседневной жизни: без них трудно представить коммерческие взаимоотношения горожан, заключение брака или сватовство. Однако предметом нашего внимания станет именно судебная клятва, которая часто замещала недостающие доказательства, служа своего рода эквивалентом истины, наряду с «судом Божьим», рассмотренным ранее.

В отдельных статьях «Барселонских обычаев», регулирующих принесение судебной клятвы, составители цитировали «Вестготскую правду», однако передавали соответствующие положения не дословно, в сокращённом или переработанном виде\footnote{\emph{Попова}\emph{ }\emph{Г}\emph{.}\emph{А}\emph{. }Указ. соч. С.~219.}. Это подтверждает знакомство компиляторов с вестготской правовой традицией, но не позволяет определить, пользовались ли они конкретной латинской рукописью, цитировали текст по памяти или обращались к промежуточной компиляции. Существование старокаталанских версий «Книги приговоров» свидетельствует о передаче его содержания на вернакуляре, однако текстуальная зависимость от них данных статей «Барселонских обычаев» не установлена.

Эту картину дополняют установления о клятве, включённые Ж.~Бастардесом-и-Парерой в приложение B к изданию «Барселонских обычаев». Как показал Р. д’Абадаль-и-Виньялес, большинство из них также восходит к более ранним текстам: обычай об ответственности клятвопреступников --- к капитуляриям Карла Великого, а обычай о клятве свидетелей --- к «Книге приговоров» \footnote{Usatges de Barcelona… P.~169--170.}. Таким образом, нормы о судебной клятве в «Барселонских обычаях» формировались посредством отбора и переработки установлений различного происхождения, включённых составителями в единый порядок судебной процедуры.

Даже в наиболее формализованной и лингвистически устойчивой части правового текста --- формулах --- можно проследить следы проникновения народного старокаталанского наречия в латинский язык «Барселонских обычаев». Это было связано с тем, что данные формулы фиксировали не только сам процесс принесения различных видов клятв, но и устойчивые словесные модели устного происхождения, в соответствии с которыми они должны были произноситься.

Составители «Барселонских обычаев», как правило, ограничивались описанием обязательной формулы и ритуального контекста принесения клятвы или присяги, тогда как полный текст самой клятвы в большинстве случаев в письменной фиксации отсутствовал. В наиболее формализованной части юридических формул мы встречаем следы проникновения народного старокаталанского наречия в латинский текст обычаев, поскольку они содержат подробные описания не только процесса принесения различного рода клятв, но и стандартных формул, в соответствии с которыми должна была произноситься судебная клятва. Как показала И.И.~Варьяш, это явление было характерно не только для христианской правовой традиции Пиренейского полуострова, но и для мусульманских судебников, созданных на территории Арагонской короны\footnote{\emph{Варьяш}\emph{ И.И. }Клятва пиренейских сарацин… С.~185.}.

\subsection{Терминология и процессуальная семантика клятвы: jurare и sagrament}

Латинская традиция различала два основных термина для обозначения клятвы: sacramentum и iurare. Однако в каталонских обычаях эти термины используются не просто как синонимы, а как обозначения различных аспектов единого процесса. Sacramentum обозначает клятву как сакральный акт, связанный с обращением к Богу и святым, тогда как iurare подчеркивает момент произнесения вербальной формулы.

Рассмотрим латинский текст 46 (49) статьи «Барселонских обычаев»: \emph{«Клятвой всегда должно клясться на освящённом алтаре или на святых Евангелиях; и тот, кто клянётся, в каждой клятве должен призвать в "свидетели", за исключением случаев лжи и предательства, "Бога и его святых"»}\footnote{Usatges de Barcelona / ed. J. Bastardas i Parera. Barcelona: Fundació Noguera, 1991. Us. 46 (49): «SACRAMENTVM sit omni tempore iuratum sy. per altare consecratum uel super sanctum euange. lium; et ille qui iurauerit, in omni sacramento debet mittere suo «sciente», excepto in bauzia et in tradicione, et “per Deum et hec sancta”». P.~86--87.}.

В этом установлении «Барселонских обычаев» встречается четыре термина, связанных с принесением клятвы (лат. --- sacramentum, iuratum sit, sacramento, iurauerit), что является лексически избыточным (лат. \emph{«Клятвой }\emph{дóлжно}\emph{ клясться»}), поскольку значение каждого из этих терминов связано с принесением клятвы. Это свидетельствует о том, что постепенно они начинают восприниматься как единая синтагма и сливаются воедино. Составители «Барселонских обычаев», по-видимому, не считали эти термины взаимозаменяемыми, в отличие от составителей «Вестготской правды», и не считали текст перегруженным повторами, поэтому они использовали слова с разными корнями в разных смыслах, вербально разграничивая феномен и процесс, действие.

Г.А.~Попова подчеркивала, что в вестготской правовой традиции наблюдалось постепенное вытеснение термина sacramentum термином iuramentum, что в дальнейшем закрепилось в старокастильском языке\footnote{\emph{Попова}\emph{ }\emph{Г}\emph{.}\emph{А}\emph{. }Указ. соч. С.~220.}. При этом в каталонской правовой традиции как связная синтагма «sacramentum sit <…> iuratum» (букв. \emph{«Клятвой }\emph{дóлжно}\emph{ <…> клясться»}) встречается лишь в тех статьях «Барселонских обычаев», которые описывали процесс принесения судебной клятвы; присяга же (лат. --- \emph{sacramentum}\emph{ }\emph{fidelitatis}), включённая в ритуал оммажа, обычно обозначалась лишь одним из терминов\footnote{Usatges de Barcelona… Us. 47 (50): «OMNES HOMINES, tam milites quam rustici, iurent senioribus suis...». P.~86--87.}.

Такой способ передачи смысла, вероятно, связан с особым процессуальным значением данной синтагмы. Х.~Мартинес-Гаскес полагает, что подобные «плеоназмы» в тексте «Барселонских обычаев» объясняются не пробелами в знании письменного латинского языка у составителей, а их стремлением придать этому правовому памятнику больший вес, наполнив его чинными повторениями\footnote{\emph{Martínez Gázquez J.} El latín en los Usatges de Barcelona // Glossae: European Journal of Legal History. 1995. Vol.~7. P.~108.}. На наш взгляд, легисты-компиляторы довольно чётко разграничивали «клятву» как явление правовой культуры и процесс её осуществления, поскольку могли часто сталкиваться с клятвами и присягами в различных контекстах.

В «Обычаях Тортосы», записанных на старокаталанском языке во второй половине XIII~в., принесение клятвы чаще всего обозначается конструкцией «принести клятву» (старокат. --- \emph{fer}\emph{ }\emph{sagrament}) или, значительно реже, глаголом «клясться» (старокат. --- \emph{jurar}). Существительное «клятва» (старокат. --- \emph{jurament}), известное по кастильской правовой традиции, в тексте памятника вовсе не употребляется; в отдельных случаях встречается также конструкция «дать клятву» (старокат. --- \emph{dar}\emph{ }\emph{sagrament}). Иногда эти термины встречаются уже как однородные сказуемые, что снова указывает на различие в их процессуальном значении, например, в статье 2.18.3:\emph{ «Если по приговору кому-либо }[предписано принести--А.К.]\emph{ клятву, или она приносится одной стороной другой стороне, и тот человек, которому [другой --А.К.] должен поклясться }(кат. --- \emph{jurar}\emph{) и клятву совершить }(кат. ---\emph{ }\emph{fer}\emph{ }\emph{sagrament}\emph{), отказывается от неё либо освобождает другую сторону от обязанности её принести, это имеет ту же силу, как если бы клятва была принесена, и впредь по данному делу не может и не должно быть предъявлено никакого иска»}\footnote{Costums de Tortosa. 2.18.3: «Si sagrament ésjugiat per sentència a alcun, o és dat de part a part, e aquel a qui hom deu jurar e fer lo sagrament, lo y lexa o lo y perdona, tot aytant val com si jurat avia, e jamés d' aquela cosa no pot ne deu ésser moguda demanda…». P.~128.}.

Возможно, эта тенденция созвучна той, которую зафиксировала И.И.~Варьяш на материалах более поздних текстов мусульманских судебников на территории Арагонской Короны: \emph{«И все же остаётся впечатление, что кастильская "}\emph{jura}\emph{" подразумевала исключительно сакральную формулу, в то время как каталанский "}\emph{sagrament}\emph{" включал в себя сакрально-юридическую формулу, объединявшую обращение к Богу, призыв к нему как к свидетелю и суть дела»}\footnote{\emph{ }\emph{Варьяш}\emph{ И.И. }Клятва пиренейских сарацин… С.~185.}. Однако по отношению к нашим источникам гипотеза, предложенная Ириной Игоревной, нуждается в некотором уточнении.

В данном случае различие носило процессуальный характер: глагол «клясться» (старокат. --- \emph{jurar}) обозначал произнесение установленной словесной формулы, тогда как конструкция «принести клятву» (старокат. --- \emph{fer}\emph{ }\emph{sagrament}) охватывала более широкую совокупность слов и сопровождавших их действий, благодаря чему клятва приобретала перформативный характер.\footnote{\emph{Остин Дж.} Перформативы -- констативы // Философия языка / ред.-сост. Дж. Р. Сёрн; пер. с англ. Изд. 2-е. М.: Едиториал УРСС, 2010. С.~23.}.

Другие варианты глаголов, обозначающих принесение клятвы, в «Обычаях Тортосы» встречаются довольно часто. В контексте разрешения тяжб можно найти следующие словосочетания: «дать клятву» (старокат. --- \emph{donar}\emph{ }\emph{sagrament}), «принять клятву» (старокат. --- \emph{pendre}\emph{ }\emph{sagrament}), «возвратить клятву» (старокат. --- \emph{tornar}\emph{ }\emph{sagrament}) и аналогичные конструкции как в действительном, так и в страдательном залоге (например, статьи 2.18.4, 2.18.10 и др.)\footnote{Costums de Tortosa. 2.18.4: «Sagrament que és donat de part a part, si aquel a qui és donat no•l vol fer o no·l vol tomar o referre a l'altra part, o serà donat per júgie, e aquel a qui seràjugiat o donat no·l volràfer, és aüt e tengut per confés, e lo júgie pot enantar contra él axí com a confés». P.~128; Ibid. 2.18.10: «…e lo demanador dirà delant lo júgie: “vàgia enant e jur”, e lo demanat se mourà de son loch, e serà al loch, e ajenolat per fer lo sagrament...». P.~129--130.}. Таким образом, под термином sagrament могло пониматься не только сакрально-правовое действие, но и само обязательство, по отношению к которому оно совершалось.

\subsection{Устная формула клятвы и её письменная репрезентация}

Прямых указаний на письменную фиксацию текстов клятв в «Обычаях Тортосы» мы не встретили. Инструкция для судебного писца, приведённая в статье 1.4.5, гласила:\emph{ «Судебный писец должен был, после принесения клятвы, как это было установлено выше, заносить в книгу [протоколов. --- прим. А.К.] жалобы и иски, возражения и изложение позиций сторон, со всеми ответами на них, поручительствами, показаниями свидетелей и всеми прочими письменными актами, относящимися к судебному разбирательству <…>»}\footnote{Ibid. 1.4.5: «L'escrivan de la cort deu escriure e-llibre aprés que•L sagrament auràfet, axí com dit és, los clams e els libels, excepcions e posicions, ab totes Lur respostes, fennançes, lo dit dels testimonis e totes altres escriptures que a plet pertaynen...». P.~39--40.}. Упоминаемый в этой статье термин sagrament, по всей видимости, обозначала процессуальную клятву стороны (или сторон), после принесения которой начиналась официальная письменная фиксация судебного разбирательства.

Записи клятв не упоминаются и в перечне свидетельств, относящихся к делу, сведения о которых должны были вноситься писцом в протокол судебного заседания. Вопрос о том, могли ли они входить в состав \emph{«прочих письменных актов, относящихся к судебному делу»}\footnote{Ibid. 1.4.5: «…totes altres escriptures que a plet pertaynen...». P.~40.}, остаётся пока неясным, поскольку составители «Обычаев Тортосы» не оставили никаких пометок и комментариев по этому поводу, в то время как иные средства доказательства они подвергли чёткой фиксации.

Кроме того, довольно трудно представить, каким образом в средневековом суде могли бы фиксировать процесс принесения клятвы, опираясь исключительно на тексты нормативного характера без привлечения судебных протоколов. В данном случае нам придётся обратиться к научной литературе, в частности к разделу книги И.И.~Варьяш «Мусульманская Европа. Сигналы идентичности», посвящённому судебным клятвам пиренейских мусульман в позднее Средневековье (XIV--XV~вв.)\footnote{\emph{Варьяш}\emph{ }\emph{И}\emph{.}\emph{И}\emph{.} Мусульманская Европа. Сигналы идентичности… С.~101.}. Для более позднего периода достоверно известно, что в ходе судебного процесса клятвы и показания свидетелей тщательно записывались именно в таком виде, в каком произносились и на том наречии, на котором они были произнесены\footnote{\emph{Она же. }Клятва пиренейских сарацин … С.~178.}. Это было необходимым условием для того, чтобы клятва, имевшая в себе элемент сакрального, была признана действительной.

О сакральном характере клятвы и о стремлении записать её «правильно», то есть так, что она служила реальным доказательством того или иного факта, мы писали выше. Однако за этим стремлением зафиксировать «произнесение» клятвы может стоять и явление более высокого уровня, связанное не только с религиозностью средневекового мышления, но и с определёнными представлениями о праве, истине и доказательстве. Так, о роли формулировки и языка в произнесении клятвы писал в своей работе «Истина и правовые установления» М.~Фуко: \emph{«В некоторых случаях }\emph{формулировку произносили и все равно проигрывали. Но не потому, что сказали неверную фразу, как и не потому, что их уличили в обмане, но потому, что произнесли формулировку не как положено»}\footnote{\emph{Фуко М. }Указ. соч. С.~83.}.

Язык, на котором произносилась клятва, понимание формулы участниками действа, равно как и её содержание, имели существенное значение. Важен был тот факт, что она была произнесена подобающим образом с соблюдением всех процессуальных норм, за которым следили должностные лица. Во «включениях» в записи судебных заседаний того, каким образом была произнесена клятва, буквальная фиксация «сказанного» становятся для исследователя очень важными свидетельствами о том, какое место занимал язык в представлениях средневекового человека о праве.

Сходные явления можно наблюдать в грамотах на «фаршированной латыни», происходящих из Южной Франции и представляющих собой клятвы верности сеньору, то есть, по сути, феодальную присягу\footnote{\textsuperscript{ }\emph{Belmon J., Vielliard F.} Latin farci et occitan dans les actes du XIe siècle // Bibliothèque de l'École des chartes. 1997. Vol.~155, №~1. P.~168; \emph{Katsura H. }Serments, hommages et fiefs dans la seigneurie des Guilhem de Montpellier (fin XIe -- début XIIIe siècle) // Annales du midi. 1992. Vol.~104. P.~141--161.}. В них отдельные глаголы и даже целые предложения записаны на окситанском и старопровансальском языках в то время, как весь основной текст грамоты, включая формуляр, был написан на латыни. С чем могло быть связано такое внимание к языку, на котором произносилась клятва, даже если он отличался от языка остального документа?

В «Барселонских обычаях» мы не найдём таких сведений, однако формулировки образцовых клятв мы в них встречаем, как в латинской, так и в каталанской версии. Однако эти формулировки крайне лаконичны и немногословны; сам способ произнесения такой вербальной формулы представляет особый интерес для исследователя.

«Обычаи Тортосы» предлагают вариативные описания порядка принесения присяги в суде, зафиксированные в прямой речи и призванные служить моделями правового ритуала. Среди них находится и текст, который надлежало озвучить заёмщику в делах о взыскании средств: \emph{«…он обязан принести присягу в той форме, которую установят судьи, а им же надлежит обратиться к нему со словами: "Клянешься ли ты, что не имеешь средств уплатить долг ни полностью, ни частично? И что, за вычетом того, что Бог даст тебе для твоего пропитания и одежды, у тебя нет ничего, что давало бы основание отказать в выплате твоему кредитору или кредиторам, если таковых несколько и они предъявляют требование …"»}\footnote{Costums de Tortosa. 1.3.17: «...per sentència dels ciutadans, deu jurar enaxí com los júgies en alà eletz lo declaran, so és saber losjúgies li deven dir: “vós juraretz que no avetz de què pagar tot lo deute ne partida d'aquel; e con Déus vos o [darà], levat vostre menjar e vostre vestir, ne vós o podretz aver per neg una raon, que·n pagaretz vostre creedor o vostres creedors, si mols són aquels qui se'n clamen”». P.~35.}.

Ещё один вариант нормативного текста судебной клятвы с использованием прямой речи предлагает текст другой статьи «Обычаев Тортосы»: \emph{«…и тот, против кого хотят свидетельствовать, пусть скажет: «<…> и тот, против кого желают представить свидетелей, скажет: „Он заявляет, что докажет и приведёт свидетелей; я не желаю, чтобы это было доказано [ими---свидетелями.---прим. А.К.], но пусть он продолжит и принесёт клятву“, --- либо же он, то есть тот, кто желает представить свидетелей, поклянётся в том, что он намерен доказать»}\footnote{Costums de Tortosa. 2.18.5: «...e sobr·assò aquel volrà provar son enteniment e diu aquel que•l provarà, e aquel contra qui volran dar los testimonis dirà: “él diu que provarà e que darà testimonis; no vul que u prou, mas vàgia enant e jur” o él, so és aquel qui vol dar los testimonis, que él jur so que entén a provar». P.~128.}. Довольно детально прописаны реплики других участников судебного процесса, в то же время практически никогда не приводится сама «формула» клятвы в устах говорящего.

Одним из объяснений того, что эталонные или образцовые формулировки судебной клятвы не были включены в текст «Вестготской правды», имели крайне краткую форму в «Барселонских обычаях» и не приводились в «Обычаях Тортосы», может быть сакральная природа клятвы, которая не предполагала её использования «вне» контекста конкретной тяжбы или разбирательства, что могло бы нивелировать её значение. Либо текст произносимой клятвы был очень простым, всем понятным и известным, например, \emph{«Клянусь Богом и его святыми!»}. Тогда формулировки, предложенные составителями «Барселонских обычаев», могли быть практически буквальными; составители «Обычаев Тортосы» предпочли описательные формулировки.

\subsection{Клятва в поликонфессиональном судебном пространстве}

Процесс принесения клятвы в делах смешанного типа между иноверцами (прежде всего, иудеями) и христианами подробно описывается в нескольких статьях «Барселонских обычаев» и «Обычаев Тортосы»\footnote{Usatges de Barcelona. Us. 46 (49): «SACRAMENTVM sit omni tempore iuratum sy. per altare consecratum uel super sanctum euange. lium; et ille qui iurauerit, in omni sacramento debet mittere suo «sciente», excepto in bauzia et in tradicione, et “per Deum et hec sancta”». P.~86; Costums de Tortosa. 4.11.38: «Con testimoni deven fer jueus contra crestians o per xrestians, juren sobre la lig de Moysèn que• Is posa hom denant, axí com fan los xrestians, sobre quatre evvangelis…». P.~214.}. Благодаря этим сведениям вопрос об особенностях клятв \emph{more}\emph{ }\emph{iudaico} в Каталонии в целом и в Тортосе и Барселоне в частности стал одним из наиболее изученных в историографии интересующего нас вопроса\footnote{См. например: \emph{Muntané i Santiveri J.-X.} Anàlisi de l'estructura del jurament de les malediccions dels jueus catalans: usatge 171 // Revista de Dret Històric Català. 2014. Vol.~13. P.~9--48.}.

Это связано с несколькими обстоятельствами: во-первых, с довольно большой численностью иудейского населения Тортосы и её округи; во-вторых, с распространением на территории Европы в Средние века уничижительного ритуала иудейской присяги \emph{more}\emph{ }\emph{iudaico}.

Специальные условия и ритуальные действия, необходимые для принесения разного рода клятв иноверцами, фиксируются ещё в тексте «Вестготской правды», они определяют особый порядок клятвы иудеев при переходе в христианство и ещё в некоторых случаях (но за рамками христианского суда)\footnote{\emph{ }\emph{Попова}\emph{ }\emph{Г}\emph{.}\emph{А}\emph{. }Указ. соч. С.~219.}. В отношении норм принесения клятв иудеями «Барселонские обычаи» демонстрируют преемственность по отношению к «Книге приговоров»\footnote{\emph{García y García A. }Los juramentos e imprecaciones en los "Usatges" de Barcelona // Glossae: European Journal of Legal History. Instituto de Estudios Sociales, Políticos y Jurídicos. 1995. Vol.~7. P.~51--52.}.

В более поздних сводах обычного права мы уже имеем дело с полноценной судебной клятвой. Наглядным примером такой трансформации является статья 4.11.38 «Обычаев Тортосы», которая детально описывает порядок принесения судебной клятвы иудеями, подчёркивая его сходство с аналогичными действиями в христианском суде. Вместо Евангелий, на которых обычно приносилась христианская судебная клятва, в данном случае предполагается использовать Моисеев закон: \emph{«Что касается свидетельства, которое должны принести иудеи против христиан или в защиту их, то пусть клянутся они [иудеи.---А.К.] на Законе Моисеевом, который положат перед ними, как делают это христиане на четырёх Евангелиях»}\footnote{Costums de Tortosa. 4.11.38: «Con testimoni deven fer jueus contra crestians o per xrestians, juren sobre la lig de Moysèn que•Is posa hom denant, axí com fan los xrestians, sobre quatre evvangelis…». P.~214.}.

Эта же статья «Обычаев Тортосы» содержит описание ещё одного ритуального действа, которое, по мысли составителей этого законодательного компендиума, в определённых случаях должно было сопровождать присягу more iudaico: \emph{«Но если иудей участвует в деле о сокрытии имущества, возбужденном против христианина или сарацина, то иудеи приносят клятву: по делам стоимостью в пять су или менее --- на Законе Моисеевом; по делам свыше пяти су --- под проклятиями, а именно так: перед ними ставят сосуд с водой и зажжённую свечу. И зачитывают им проклятия над их головами; и там, где говорится, что они клянутся, он отвечает, что клянётся, а в тех местах, где следует отвечать „аминь“, они должны сказать „клянусь, аминь“. И по прочтении проклятий следует }\emph{погасить свечу в воде»}\footnote{Ibid.: «…Mas si jueu a afer escondiment a xrestian o a ssarayn, juren los jueus, de V sous o de V sous a avayl, sobre la lig de Moysèn; de V sous en amunt juren sobre les maledictions, so és que tenen denant éls I bacín ab aiga e una candela ensesa. E lig-los hom les maledictions sobre lo cap; e là on diu que juren, respon que jura, eanant là on deven respondre “amén”, que an a dir “jura amén”. E les malediccions leztes a a·pagar la candela en l'aiga». P.~214.}. Стоит отметить, что во второй редакции «Обычаев Тортосы» ритуально-уничижительная часть процесса принесения клятвы с зажиганием свечи и последующим её тушением в сосуде с водой была упразднена по религиозным причинам\footnote{Ibid.: «…excepto quod dicitur lbl de baclno aque et de candela accenssa In quo reprobant eam ne Deus in hoc vldeatur tempteri». P.~214.}.

Другая статья «Обычаев Тортосы» определяла, что судебные дела с участием иудеев и особенно с принесением ими клятвы следует разбирать в ином сакральном пространстве, а именно в синагоге: \emph{«И за исключением тяжб с участием иудеев, в которых, если ими должна быть принесена клятва, она должна совершаться в суде и затем --- в синагоге, и свидетельства иудейских свидетелей должны приниматься в синагоге, а всё прочее --- в суде»}\footnote{Costums de Tortosa. 7.7.20: «e exceptat plet de jueus que, si sagrament s'i deu fer per éls, deu-se jug iar en la cort e fer a la sinagoga, e pendre testimonis jueus a la sinagoga, e tot l' als en la cort». P.~366.}.

Тут же стоит заметить, что мы не видим аналогичных установлений в отношении мусульман, которым в отдельных случаях также требовалось принести клятву в христианском суде. Единственное упоминание о мусульманской клятве присутствует в статье 2.18.11 в редакции «Обычаев Тортосы» 1276~г., которая сообщала о том, что все иноверцы (и иудеи, и мусульмане) должны приносить судебную клятву вне здания суда и не должны наблюдать за процессом принесения клятвы христианами, даже если последние являются ответчиками, а иноверец истцом: \emph{«Когда христианин должен совершить клятву по иску, который предъявляет ему иудей, он должен предстать перед }\emph{вегером}\emph{ и горожанами. Иудей же обязан удалиться из суда там до тех пор, пока клятва не будет принесена, поскольку не следует клятву совершать в присутствии иудея, и не должен он }\emph{[иудей.---прим. А.К.] видеть, как клятва совершается. То же касается и мусульман»}\footnote{Costums de Tortosa. 2.18.11: «Conxrestian deu fer sagrament per demanda quejueufa contra él, deu-lo feren la cort denant lo veger, e els ciutadans. Mas lo jueu deu exir defora la cort e estar tro que-l sagrament sia fet, car dementre lo jueu hi és no·l deu fer es (sic), ne el jueu no deu veure fer aquel sagrament. [dem in sarracenis]». P.~130.}.

Позднее обе процитированные выше статьи были исключены из текста «Обычаев Тортосы», поскольку постоянно переносить место проведения судебного заседания церковные иерархи сочли неподобающей практикой, как указано в маргиналии к рукописи А. Как отметил Ф.~Дэйлидер, такое соотношение норм для дел смешанного типа в каталонских нормативных текстах XIII~в. могло объясняться более сильной религиозной нетерпимостью по отношению к иудейскому, нежели к мусульманскому населению\footnote{\emph{Daileader P. }La coutume dans un pays aux trois religions: la Catalogne, 1228--1319 // Annales du Midi: revue archéologique, historique et philologique de la France méridionale. 2006. Vol.~118, №~255. P.~384.}.

Интересен тот факт, что в описании практик совершения клятвы иноверцами крайне мало сведений о судебной клятве мусульманского населения (фактически единственное упоминание в объёмном тексте кодекса, хотя иные формы взаимодействия с мусульманской общиной регламентируются весьма подробно), судя по другим установлениям «Обычаев Тортосы», многочисленной в этом городе, так М.Т.~Феррер-Майоль привела данные налоговой переписи 1410~г. о 32 главах сарацинских домохозяйств в Тортосе и Тивеньсе\footnote{\emph{Ferrer Mallol M.T. }L'Aljama islámica de Tortosa a la Baixa Edat Mitjana // Recerca. Centre d'Estudis Històrics del Baix Ebre, 2003. Núm. 7. P.~186.}. Для более корректного сравнения, к сожалению, невозможно установить точное число иудеев, проживавших в XIII~в. в Тортосе, которая наряду с Льейдой имела одно из наиболее крупных городских иудейских сообществ. Из документов королевской курии известно, что после 1148~г. тортосским иудеям были пожалованы территории для строительства 60 домов, что означает, что иудейская община была также весьма многочисленной\footnote{\emph{Curto Homedes A. }Topografia del call jueu de Tortosa // Recerca. 1999. Núm. 3. P.~10.}.

Это может быть связано сразу с несколькими обстоятельствами: с одной стороны, с активной адаптацией мусульманской альхамы Тортосы к нормам королевского права и частыми обращениями в суд, с другой стороны, с тем, что составители данного свода обычного права имели относительно размытое представление о правовых практиках, господствовавших внутри мусульманского сообщества, пусть и многочисленного.

Стоит также отметить, что в записях обычного права с территории «Старой Каталонии» сведения о судебной клятве иноверцев вообще отсутствовали. Однако указания на особую форму судебной присяги для иудеев сохранились в «Обычаях Тортосы». В примечании издателя интересующего нас правового памятника указано, что это могло быть следствием общего снижения толерантности христианской власти к иноверцам, которое фиксируется во второй редакции «Обычаев Тортосы»\footnote{Costums de Tortosa. 2.18.11. P.~130.}.

\subsection{Изменение процессуального значения клятвы в XIII~в.}

В XIII~в. произошло изменение отношения легистов к клятве как способу доказательства. Наряду с сохранением значения клятвы как самостоятельного способа судебного доказательства, она приобретала и дополнительную функцию обязательного процессуального действия, предшествовавшего, в частности, даче свидетельских показаний.

В то же время для легистов, изначально составлявших текст кодекса, судебная клятва в системе доказательств не занимала главенствующего положения и ко 2-й половине XIII~в. уже считалась инструментом довольно архаичным и даже устаревшим, что показывают нормы относительно клятв и присяг христианского населения Тортосы и Барселоны.

Первоначально «Обычаи Тортосы» не предусматривали специального наказания за нарушение клятвы, принесённой в суде, ввиду её сакрального характера: \emph{«Никого не следует наказывать за ложную клятву или лжесвидетельство, поскольку за это наказывает Господь, но не люди»}\footnote{Ibid. 2.18.9: «Nul hom no és punit per sagramentfals ne per perjuri, car aquela pena a Déu s'esgarda, e no ha hòmens...». P.~129.}. Первую редакцию этого обычая также сопровождает латинская вставка, сама по себе представляющая интерес. Это изменённая цитата из Corpus Iuris Civilis, вероятнее всего, попавшая в текст «Обычаев Тортосы» из текста «Фуэро Валенсии»: \emph{«Пренебрежение религиозной сутью клятвы будет наказано Богом, поскольку кары за клятвопреступление, свыше, вполне достаточно»}\footnote{Furs de València / A cura de Colon G. i Garcia A. Barcelona: Editorial Barcino, 1974. Vol.~2. P.~254--255. «Ius iurandi contempla religio satis Deum habet ultorem nam sufficit pena periurii quam a Deo expectat»}. Как показала М.~дель~Кармен~Лáзаро~Гийамон, отказ власти от наказания за клятвопреступление со ссылкой на последующую Божью кару мог свидетельствовать о сближении клятвы и ордалии в процессе трактовки постклассических римских правовых норм\footnote{\emph{del Carmen Lázaro Guillamón M. }El iusiurandum iudicium delatum: entre medio de prueba y ordalía // La prueba y medios de prueba. De Roma al derecho moderno: actas del VI Congreso Iberoamericano y III Congreso Internacional de Derecho Romano, 2000. [Madrid]: Servicio de Publicaciones, 2000. P.~410.}.

В то же время во второй редакции данного кодекса приводится уточнение: \emph{«Этот обычай не применяется в отношении тех, кто ложно свидетельствовал, и тех, кто приносил клятву перед сеньорией по какому-либо делу, потому что такие должны быть наказаны, если будет доказано, что они лгали при присяге или были лжесвидетелями»}\footnote{Costums de Tortosa. 2.18.9: «...Si a alcunxrestian seràjugiat de fer sagrament, aquel a qui seràjugiat que·l·fassa, pot aver altra persona que siaxrestià que-l·fassa en loch d' él, él present serà, atorgant que sobre la sua ànima lofassa per alcuna cosa que li sia deu». P.~129.}. Эта статья «Обычаев Тортосы» не только является внутренне противоречивой, но и в определённом смысле противоречит и нормам 63(67)-й статьи «Барселонских обычаев», которая предполагает телесное либо денежное наказание, а также имущественную компенсацию ущерба со стороны лица, признанного клятвопреступником\footnote{\textsuperscript{ }Usatges de Barcelona. Us. 63 (67): «…Et postea non testificetur in placito nec credatur per sacramentum».P.~100--101.}.

Кроме того, согласно 63 статье «Барселонских обычаев», клятвопреступник не мог «\emph{ни свидетельствовать в суде, ни клятву приносить»}\footnote{Ibid.: «…Sin autem et periurus inde esse uidetur, aut manum perdat aut centum solidis redimat aut quartam partem facultatum suarum amittat, profuturam in manu illius cuius periurus effectus est…». P.~100--101.}. Наказания за клятвопреступление, а именно отрубание руки и уплата 100 солидов, имели явно франкское происхождение и не совпадали с нормами «Вестготской правды».

Судебная клятва в каталонском обычном праве XII--XIII~вв. служит индикатором глубокого перехода от устной к письменной правовой культуре. Сочетание устного ритуального произнесения с письменной фиксацией в судебных записях отражало переходный характер этого периода. Судебная клятва постепенно теряла свою доказательную функцию, но при этом приобретала новое значение как элемент процедуры, символизирующий переход к более формализованной и письменно-ориентированной правовой культуре.

Ни в «Барселонских обычаях», ни в «Обычаях Тортосы» клятва не выступала универсальным или приоритетным средством доказательства: её применение зависело от характера дела, наличия свидетелей и иных доказательств, а также от процессуального выбора сторон. При этом сопоставление двух сводов показало, что в «Барселонских обычаях» клятва была тесно связана с ранними традиционными формами доказательства: ордалиями и судебными поединками, тогда как в «Обычаях Тортосы» она всё чаще функционировала как элемент процедуры, встроенный в более рационализированную систему свидетельских показаний и письменной фиксации. Эволюция терминологии, связанной с принесением клятвы, отражает это смещение: различие между сакральным актом и процессуальным действием постепенно размывается, а клятва утрачивает статус автономного доказательства, оставаясь значимым перформативом.

Сопоставление норм о клятвах христиан, иудеев и мусульман позволяет сделать вывод о том, что конфессиональные различия оказывали существенное влияние на процессуальные формы, но не разрушали единства судебной системы. Особые ритуалы, ограничения пространства и присутствия сторон, а также вариативность формулы клятвы more iudaico свидетельствуют о попытке христианской власти вписать их в общий судебный порядок посредством специально оговорённых процедур. В более широком контексте судебная клятва предстает индикатором переходного состояния каталонского права: от устной, ритуально насыщенной культуры доказательства к институционально оформленной судебной практике. Именно эта двойственность --- сочетание устного и ритуального характера клятвы с развивающимся процессуальным правом под влиянием рецепции римского права --- определила место клятвы в правовой культуре Каталонии XII--XIII~вв. и, по крайней мере частично, может объяснить её постепенную трансформацию из средства доказательства в формализованный элемент судебного процесса.

\section{Репутация как основание процессуального доверия}

Особое место в системе судебного доказательства в Каталонии XII--XIII~вв. занимали репутационные способы доказательства --- прежде всего представления о доброй славе, общественном мнении, «цельной» или «безупречной» репутации лица. В более широком европейском контексте средневековая «слава» (лат. --- \emph{fama}) обозначала одновременно и «то, что говорится», и возникающую из этого коллективного говорения репутацию\footnote{\emph{Bowman J.A.} Infamy and Proof in Medieval Spain // Fama: The Politics of Talk and Reputation in Medieval Europe / ed. Fenster T., Lord Smail D. Cornell University Press, 2019. P.~95; \emph{Théry}\emph{~}\emph{J. }“Fama”: Public Opinion as a Legal Category. Inquisitorial Procedure and the Medieval Revolution in Government (12th--14th centuries) // Micrologus. 2024. Vol.~32. P.~153--193.}; уже с конца XII~в. она все отчетливее превращалась в юридически значимую категорию, связывающую суд, общественное знание и публичную оценку личности. Иначе говоря, в данном случае мы имеем дело с признанной общиной способностью быть субъектом права: обвинять, свидетельствовать, клясться, защищать свою честь и доверять своему слову.

\subsection{Статус, честь и доверие к участнику процесса}

В «Барселонских обычаях» термин добрая слава в явном виде, по-видимому, не занимает столь заметного места, как в более поздних юридических памятниках, например, «Обычаях Тортосы». Однако сама идея репутационной допустимости здесь присутствует отчетливо. Так, клятвопреступникам нельзя было выступать в качестве свидетелей, поскольку их репутация уже была небезупречной: \emph{«Если же выяснится, что он повинен в клятвопреступлении <…> В дальнейшем такой человек не вправе ни свидетельствовать на суде, ни приносить судебную клятву»}\footnote{Usatges de Barcelona. Us. 63 (67): «Sin autem et periurus inde esse uidetur <…> Et postea non testificetur in placito nec credatur per sacramentum». P.~100.}\emph{.} Запрет клятвопреступникам участвовать в процессуальных действиях показывает, что утраченное доверие к лицу влечет утрату его дееспособности как свидетеля: человек, однажды нарушивший клятву, более не может ни заслуживать веры, ни подкреплять ею свои показания в качестве свидетеля. Тем самым репутация выступает не дополнением к доказательству, а его предварительным условием.

Таким образом, в «Барселонских обычаях» клятва, статус, честь и свидетельская пригодность еще не вполне выделились в строго разграниченные юридические категории. Для этого памятника характерна их тесная связь: возможность свидетельствовать зависела не только от знания фактов, но и от общественной оценки, а клятва подтверждала не только позицию стороны по делу, но и ее нравственную и социальную состоятельность. В условиях, когда письменная фиксация прав оставалась неполной и не всегда доступной, добрая слава выполняла важную доказательную функцию, отчасти восполняя недостаток документа. Репутация служила признаком того, что показаниям данного лица можно доверять. Поэтому в «Барселонских обычаях» она удостоверяла не отдельное обстоятельство дела, сколько процессуальную надежность самого участника разбирательства.

В то же время, например, «Обычаи Льейды» также не выделяли отдельных категорий «доброй» и «дурной славы», ограничиваясь лишь термином «подходящего» по отношению к свидетелям, который можно перевести также как «достойных доверия»: «\emph{Также установлено, что они не вправе ни сами, ни через своих }\emph{байлов}\emph{ предъявлять нам обвинение или упрек в чем бы то ни было без пригодных (надлежащих) свидетелей»}\footnote{Costums de Lleida. [5]: «De inculpatione nobis non fatienda. Item, quod non possint per se vel per baiulos suos nos inculpare aut increpare de aliquo absque testibus idoneis». P.~99.}. Иных упоминаний о доброй или дурной славе в данном кодексе нами найдено не было. По нашему мнению, это не говорит о том, что эта правовая категория отсутствовала в правовой культуре Льейды, скорее, она просто не была подробно отрефлексирована в данном кодексе.

XX статья «Кутюмов Перпиньяна» еще отчетливее показывает, что судопроизводство в этом городе было тесно связано с защитой социальной иерархии его обитателей: если \emph{«низкий человек или }\emph{бакалларий}\emph{»} нанесет оскорбление \emph{«доброму человеку Перпиньяна»}, то любой присутствующий может немедленно одернуть его во время самой ссоры, не причиняя телесного вреда, и сеньор не вправе требовать ничего с того, кто это сделал\footnote{Les Coutumes de Perpignan. [XX]~: «Item, si vilis persona vel bacallator iniuriam fecerit vel dixerit alicui probo homini de Perpinyano, alius circumstans potest eum corripere in ipsa ricxa sine deterioratione illius persone; et quod dominus nichil potest petere ab eo qui eum corripuit». P.~573.}. Следовательно, защита общественного порядка здесь частично выводилась за пределы формального суда и возлагалась на самих горожан, поскольку община получала право на непосредственное пресечение поведения, нарушающего признанную статусную иерархию.

Иное, более развитое состояние категории доброй или дурной славы обнаруживают «Обычаи Тортосы». Именно в Тортосе добрая и дурная слава, а также близкая ей «безупречная слава» (старокат. --- \emph{d’íntegra}\emph{ }\emph{fama}), стали заметными процессуальными категориями, которые в каждой редакции упоминаются не менее 20 раз. Такая частота встречаемости данных терминов в «Обычаях Тортосы» свидетельствует о том, что это устойчивая лексика правового отбора. Если для «Барселонских обычаев» можно говорить преимущественно о репутации как предпосылке доверия, то в Тортосе она уже систематически встроена в судебный механизм: регулирует право на обвинение, применение пытки, статус свидетельства и даже ход розыскного процесса.

\subsection{Добрая слава и условия участия в судебном процессе}

Показательно прежде всего правило о допустимости обвинителя: «\emph{Кто хочет обвинить другого <…> не должен быть преступным человеком, но человеком доброй славы»}\footnote{Costums de Tortosa. 9.1.3: «Qui accusar vol altre <…> no deu ésser criminós, mas hom de bona fama». P.~408.}. Право на обвинение предполагало еще и социальную пригодность обвинителя. Эта норма показывает, что для обращения в суд необходимо было иметь определённую репутацию внутри городской общины, что означало юридическое признание социального капитала.

Здесь особенно ясно видно, что добрая слава в Тортосе не подменяла возможность обвинения, а лишь открывала к нему доступ. Она делает возможным обвинение, тогда как ее отсутствие блокирует саму инициативу процесса. В таком устройстве правопорядок оказывается глубоко коммуникативным: речь о преступлении не принадлежит никому в отдельности, а должна исходить из «правильного» социального места. Обвинение, произнесенное лицом дурной славы, с самого начала утрачивает часть своей юридической силы, потому что общество не готово признать за таким человеком право говорить от имени справедливости.

Не менее показателен текст о пытке, который определяет социальную принадлежность тех, кто не мог обладать доброй славой: \emph{«Никто, кто пользуется доброй славой, не может и не должен быть подвергнут пытке. }\emph{Но люди дурной славы и низкие лица --- как, например, }\emph{бастайши}\footnote{\textbf{Бастайши}\textbf{ }(кат. bastaixos; вероятно, от вульг.-лат. bastaxus < греч. bastázō) --- носильщики, профессиональные переносчики тяжестей. Cм. Bastaix // diccionari.cat. URL: \url{https://www.diccionari.cat/GDLC/bastaix\#} (дата обращения: 14.03.2026).}\emph{, пьяницы, завсегдатаи таверн, рабы, мошенники, игроки, ростовщики на игре <…>, }\emph{алкавоты}\footnote{\textbf{Алкавот}\textbf{ }(кат. alcavot, мн. ч. alcavots; от араб. al-qawwād --- «сводник, посредник в любовной связи») --- сводник, посредник в внебрачных или иных «незаконных» любовных связях. Cм. Alcavot // Diccionari de la llengua catalana. URL: \url{https://dlc.iec.cat/Results?EntradaText=alcavot&OperEntrada=0} (дата обращения: 14.03.2026).}\emph{ и им подобные, --- [могут.---А.К.]»}\footnote{Costums de Tortosa. 9.5.2~: «Negun no pot ne deu ésser posat a turment, que sia de bona fama. Mas hòmens de mala sia de bona fama e que sien vils persones, axí con són bastaix, e bevedors, embriachs, e tavents, o servus, bevedors, hòmens trixadors, jugadors, e prestadors de joc<...> e alcavotes, e altres semblantz». P.~422.}. Эта норма чрезвычайно важна, поскольку показывает: репутация действует не только как привилегия доверия, но и как привилегия телесной неприкосновенности. Добрая слава выступала защитой от крайней формы процессуального принуждения, тогда как дурная слава, напротив, превращала тело в допустимый объект насильственного извлечения истины.

Список лиц низкого социального происхождения, приведённый в этой статье «Обычаев Тортосы», также заслуживает внимания. Перед нами не нейтральное перечисление профессий или привычек, в один ряд поставлены низкий трудовой статус, зависимость, пьянство, азартная игра, сексуальная нечистота и мошенничество. Таким образом, воспроизводится моральная иерархия городской общины. Правовая недостоверность здесь возникает на пересечении этического, сословного и поведенческого критериев. Право исходит из представления, что некоторые формы жизни сами по себе делают показания человека «сомнительными». В репутационном доказательстве устанавливается не только истина по делу, но и граница между социально признанным и социально маргинальным.

Ещё отчётливее репутационная логика видна в составе норм о свидетелях. Статья «Обычаев Тортосы» 9.25.14 требует, чтобы свидетелями были \emph{«люди безупречной и доброй славы, не враги и не те лица, которых }\emph{право устраняет из числа свидетелей»}\footnote{Costums de Tortosa. 9.25.14: «Los testimonis deuen ésser aquels qui sien d’entrega fama e de bona, e no enemichs…». P.~490.}. Далее они должны клясться, что не скажут лжи в том, что обязуются говорить только правду о том, что знают или видели. Здесь особенно существенна двойная фильтрация: с одной стороны, проверяется репутация свидетеля; с другой --- его возможная заинтересованность. Но и то и другое относится не к содержанию будущих показаний, а к самому социальному положению говорящего. Право заранее определяет, какие виды связи с общиной делают лицо пригодным к свидетельству, а какие --- напротив.

В 9.25.20 тот же принцип получает процедурное развитие: \emph{«Если обвиняемый отрицает предъявленные ему пункты, }\emph{вегер}\emph{ по решению инквизиторов вызывает свидетелей доброй славы, которые не являются его врагами и вообще не подпадают под подозрение; в присутствии обвиняемого они приносят клятву говорить правду, затем их показания записываются и оглашаются, а обвиняемому или его представителю дается возможность ознакомиться с именами свидетелей и содержанием показаний и защищаться самому либо через других лиц»} \footnote{Ibid. 9.25.20: «E los enqueridors ab lo veger present, reeben sa confessió e meten-la en escrit. E si l'encolpat, los capítols contra él posatz, totz o part, negarà, lo veger, per juydels enqueridors, cit los testimonis que sien de bona fama e no sien enemics seus ne d'altrement sopitosos...». P.~493.}. Даже в рамках розыскного процесса репутация свидетеля остается основой допустимости его слова. Следовательно, в Тортосе рационализация процесса не устраняет репутационный критерий, а переводит его в более формализованные формы: добрая слава становится частью институциональной процедуры отбора доказательства.

Более поздние сборники обычаев Орты (1296) и Миравета (1319), составители которых основывались на материалах «Льейды», также активно использовали институт «доброй славы» как одну из гарантий достоверности свидетельских показаний. В «Обычаях Орты» содержится следующая норма: \emph{«Пусть рассудительные женщины доброй }\emph{славы могут свидетельствовать в гражданских делах так же, как мужчины»}\footnote{Costums d’Orta. [XLII]: «Item, quod mulieres discrete et bone fame possint facere testimonium in causis civilibus prout mares». P.~622.}.\emph{ }Требование, чтобы женщина-свидетельница была «рассудительной» и «доброй славы», следует понимать в контексте общего недоверия средневековой правовой традиции к женскому публичному слову. Под влиянием римско-правовых и канонических представлений женщина мыслилась как существо более «слабое», подверженное внешнему влиянию и потому менее пригодное к участию в публичной судебной сфере\footnote{\emph{Kelleher M.A.} The Measure of Woman: Law and Female Identity in the Crown of Aragon. University of Pennsylvania Press, 2010. P.~26--27.}. Однако это недоверие не исключало женского свидетельства полностью: напротив, в тех делах, где женщины обладали специфическим знанием фактов, их показания признавались допустимыми. На этом фоне норма «Обычаев Орты» выглядит как компромиссная и вместе с тем показательная: женское свидетельство допускается лишь при условии, что оно подкреплено социальной репутацией и способностью к благоразумному суждению. Тем самым право конструирует доверие к женщине-свидетельнице не как к нейтральному носителю факта, а как к лицу, чья процессуальная пригодность должна быть специально удостоверена.

При этом 133 статья «Обычаев Льейды», допускавшая женщин к свидетельству только по делам о происшедшем в женских банях, отражает общее для средневекового права представление о женском слове как ограниченно пригодном для публичного судебного пространства. В «Обычаях Миравета» эта же модель получила уточнение: даже в тех случаях, когда женское свидетельство признается возможным --- по делам, связанным с баней (ст. 105), --- оно могло приниматься лишь от женщин доброй славы\footnote{Constitutiones baiulie Mirabeti. 105: «Femine bone fame recipiantur et peribeant testimonium in facto fumi et balnei; sset in aliquo [alio] caso non valeat earum testimonium». P.~27.}. Тем самым право связывает допустимость женского слова не только с его предметной релевантностью, но и с репутационной благонадежностью свидетельницы, укорененной в коллективном знании общины о ее поведении. Особенно наглядно это проявилось в статье 113 о внебрачном ребенке: \emph{«Если женщина, утверждающая отцовство определенного мужчины, известна как состоявшая в связи только с ним, ее присяга влечет обязанность мужчины содержать ребенка; если же за ней закрепилась репутация женщины, вступающей в связь и с другими, ее слова уже недостаточно»}\footnote{Ibid. 113: «...Set si dicta mulier fuerit fame vel oppinionis quod ipsa utatur cum illo et cum aliis, ille non cogatur nutrire illam creaturam~». P.~30.}. Таким образом, в Миравете женские свидетельские показания и женская клятва получали юридическую силу лишь тогда, когда они опирались на добрую славу. Репутация здесь выступала не только нравственной оценкой, но и формой коллективного знания общины.

Категория доброй или дурной славы могла относиться не только к женщинам: так, 104-я статья «Обычаев Миравета» гласит: \emph{«Если истец --- человек доброй славы, то до суммы в 20 солидов он может доказывать свое требование одним достойным доверия свидетелем либо несколькими свидетелями…»}\footnote{Ibid. 113: «...si aliquis petitor erit bone fame… valeat… probare cum uno teste fidedigno…». P.~28.}. Таким образом, слово уважаемого человека в суде получает больше внимания.

Пример «Обычаев Миравета» показателен тем, что репутация здесь выступает не только мерой публичной чести, но и формой коллективного знания общины о повседневном поведении человека. Суд исходит не столько из «внутренней» истины частной жизни, сколько из того, как эта жизнь воспринимается и оценивается окружением. Поэтому в делах о происхождении ребенка, домашнем быте и малозначительных денежных требованиях добрая или дурная слава фактически становится сокращённым видом доказательства: она позволяла восполнить недостаток формализованных средств установления факта. В этом отношении «Обычаи Миравета» развивают линию, уже заметную в «Обычаях Льейды», где женское свидетельство допускалось лишь в строго ограниченных случаях, и вместе с тем сближаются с «Обычаями Тортосы», в которых добрая и дурная слава\emph{ }получили отчетливое юридическое выражение. По сравнению же с «Барселонскими обычаями», где репутационный критерий еще не оформлен столь последовательно и проявляется главным образом через ограничения процессуальной дееспособности, в сводах обычного права более позднего периода (2-я пол. XIII-- начало XIV~в.) эти репутационные категории превращаются в самостоятельный юридический ресурс, посредством которого суд использует социальную видимость и коллективную память общины для установления истины.

\subsection{Дурная слава: подозрение, ограничения и восстановление репутации}

Противоположная доброй славе -- слава дурная в «Обычаях Тортосы» выступает не только бытовая характеристика, но и является чёткой процессуальной категорией, определявшая допустимость и доказательную силу свидетельского слова. Это особенно ясно видно в нормах о свидетелях. Так, статья 4.11.8 «Обычаев Тортосы» устанавливает, что для установления обстоятельств дела достаточно двух или более свидетелей,\emph{ «если только они не дурной славы»}\footnote{Costums de Tortosa. 4.11.8: «...ab que no sien de mala fama...». P.~208.}\emph{. }Наряду с этим ту же функцию могут выполнять публичная грамота, завещание или иное письменное распоряжение, подтвержденное двумя или более свидетелями. Дурная слава действовала здесь как отрицательный критерий процессуальной пригодности: свидетель формально мог существовать, но если он обладал дурной славой, то его показания утрачивали должную юридическую ценность.

Тот же принцип закреплялся и в статье 4.11.11: показания с чужих слов не имели силы, а одного свидетеля, даже если он присутствовал при всем происшедшем, было недостаточно для вынесения решения; исключение делается лишь для исков до 50 су, когда допускается показание одного свидетеля, но и в этом случае при непременном условии, что и сам свидетель, приносящий присягу, является \emph{«честными и не дурной славы»}\footnote{Ibid. 4.11.11: «...honest e no de mala fama. P.~208--209.}. Иначе говоря, репутация свидетельствующего лица здесь встроена в саму структуру судебного разбирательства: она не заменяла собственно доказательство, но служила условием его действительности.

Особенно показательно, что дурная слава в «Обычаях Тортосы» понимается не как абсолютное и неизменное качество человека, а как характеристика, значимая в конкретный момент. В~статье~4.11.18 указано, что свидетели, присутствовавшие при составлении завещаний, грамот, тяжб или иных договоров, должны пользоваться доверием, «\emph{если в то время, когда они давали свидетельство или были вписаны в документ, они не были дурной славы,} \emph{даже если позднее они сделались бесчестными или «дурной жизни», их прежнее свидетельство сохраняет силу»}\footnote{Costums de Tortosa. 4.11.18: «…si ladones con fan lo testimoni o són escritz <...>no són de mala fama, deven ésser creegutz axí com a bons e a leyals, jasia so que depux se sien fetz infames o de mala vida. P.~210.}. В данном случае право интересует не абстрактная нравственная оценка личности, а ее общественная репутация в конкретный момент процессуально значимого действия. Соответственно, дурная слава --- это скорее правовая, а не моральная категория, чётко привязанная ко времени. Этим объясняется и возрастной ценз, закрепленный в статье 4.11.19: «\emph{в уголовных делах допускаются свидетели старше двадцати лет, но лишь при условии, что они не дурной славы»}\footnote{Ibid. 4.11.19: « …Mas si són majors de XX ayns conplitz, són e deven ésser reebutz, ah que no sien de mala fama en los criminals pleitz. ab que no sien de mala fama en los criminals pleitz». P.~210.}. Следовательно, способность свидетельствовать по уголовным делам определяется не только возрастом и формальной дееспособностью, но и репутационной полноценностью лица. Суд, таким образом, отбирает не просто тех, кто что-то видел или слышал, а тех, чье слово признается социально достойным доверия.

Однако значение дурной славы в «Обычаях Тортосы» не исчерпывалось индивидуальным уровнем. В рамках розыскного процесса это понятие переносится с отдельного лица на городскую общину в целом. Именно так можно понимать начало статьи 9.25.1, где говорится: \emph{«…Если совершаются деяния против Бога и против справедливости, по землям распространяется дурная слава…»}\footnote{Ibid. 9.25.1: «…E Jos contra Déu e contra justícia, e se'n seguís mala fama per les terres…» P.~487.}, то даже при отсутствии прямой вины самих граждан Тортосы по соглашению между горожанами и сеньорией должно быть проведено расследование скрытых преступлений, \emph{«дабы в указанном городе могла быть осуществлена справедливость»}\footnote{Ibid. 9.25.1: «…per so que justícia pogés ésser seguida en la dita ciutat sobre Vlll capítols dejús escritz e nomenatz…». P.~487.}.

В этом контексте \emph{«дурная слава» }относится уже не к свидетелю, обвинителю или обвиняемому, а к самой Тортосе как политико-правовому субъекту. Нераскрытые злодеяния порождают не только пробел в знании, но и репутационный ущерб: город начинает пользоваться дурной славой в окружающих землях, а значит, компрометируется как пространство, в котором справедливость не действует должным образом. Поэтому \emph{розыскной процесс }вводится здесь не просто как средство выявления скрытого преступления, но и как механизм восстановления коллективной чести города и доверия к его юстиции. Иными словами, в «Обычаях Тортосы» \emph{дурная слава} оказывается двоякой категорией: на индивидуальном уровне она ограничивает процессуальную пригодность лица, а на уровне городской общины становится признаком институционального неблагополучия, требующего официального расследования.

\subsection{Индивидуальная и коллективная репутация}

В рассмотренных каталонских источниках не зафиксировано классической компургации --- коллективной клятвы группы лиц, подтверждающих правоту стороны\footnote{Подробнее о практике компургации на материалах другого региона см. \emph{Винокурова}\emph{~}\emph{М.В.} Обвинение, оправдание, наказание в обычном праве малых городов средневековой Англии // Новое прошлое / The New Past. 2016. №~3. С.~117--128.}. Однако это не означает исчезновения коллективного начала из судебного доказательства. Оно сохраняется в иных формах: в требовании доброй славы обвинителя и свидетелей, в исключении лиц дурной славы, в учете вражды, подозрительности и вообще в опоре суда на знание общины о человеке. Иначе говоря, коллективное подтверждение правоты не исчезает, а переносится из сферы совместной клятвы в сферу свидетельства и репутации.

Поэтому категории доброй и дурной славы следует понимать не как побочный или архаический элемент процессуального права, а как один из его основных механизмов для рассматриваемого нами периода. В условиях, когда документ был доступен не всегда, архаичные формы доказательства (ордалии и судебные поединки) постепенно уходили в прошлое, а следственные формы розыскного процесса только складывались, суд опирался на коллективно признаваемую шкалу доверия: кто вправе обвинять, свидетельствовать, приносить клятву, чье слово заслуживает веры, а чье --- нет. Репутация тем самым превращалась в юридический ресурс и форму социальной памяти.

Различие между «Барселонскими обычаями», «Обычаями Льейды» и «Обычаями Тортосы» состоит не в простом противопоставлении «архаичного» и «рационального», а в степени юридической проработанности репутационного критерия. В «Барселонских обычаях» он был связан с моральной и статусной пригодностью участника процесса. В «Обычаях Льейды» и созданных на их основе «Обычаях Орты» и «Обычаях Миравета» добрая и дурная слава уже отчетливо влияли на допустимость свидетельства. В «Обычаях Тортосы» же репутационные категории получили наиболее развитое процессуальное значение, определяя право обвинения, пригодность свидетеля и даже основания для розыскного расследования. Таким образом, в каталонском городском праве XII--XIII~вв. истина во многом оставалась социальной истиной: право превращало мнение общины о человеке в дополнительную форму доказательства.

\section{Свидетельское показание и оценка его достоверности}

В системе судебного доказательства допрос и свидетельские показания занимали центральное место, поскольку именно они обеспечивали суд сведениями о спорном факте там, где письменные доказательства отсутствовали, были недостаточны или не считались исчерпывающими. Вместе с тем в каталонском обычном праве XII--XIII~вв. свидетельство понималось не как абстрактное сообщение о факте, а как процессуально и социально обусловленное высказывание. Значение имели не только содержание показаний, но и личность свидетеля, обвинителя или обвиняемого: его репутация, возраст, пол\footnote{Об этом см. §3.5 «Репутационные формы установления истины» настоящего исследования.}, правовой статус, отношение к сторонам. В силу этого нормы о допросе и свидетельских показаниях позволяют проследить, каким образом суд соединял установление факта с оценкой социальной достоверности говорящего.

Сопоставление «Барселонских обычаев», «Обычаев Льейды» и «Обычаев Тортосы» позволяет проследить, как в каталонском праве XII--XIII~вв. менялись критерии допустимости и формы допроса как процессуального действия и сама оценка свидетельского слова в суде.

\subsection{Условия допуска к свидетельству}

Если исключить клятву, ордалии и судебный поединок, «Барселонские обычаи» позволяют реконструировать достаточно оформленную модель судебного разбирательства, основанную на последовательности процессуальных действий, участии сторон и судей, а также на использовании свидетельских показаний и письменных актов. 21-я статья «Барселонских обычаев» устанавливала сроки вызова в суд в зависимости от статуса участников: \emph{«Судебное заседание должно назначаться магнатам и рыцарям сначала за десять дней, а затем --- через каждые восемь дней; }\emph{крестьянам же оно должно назначаться на четвертый или пятый день»}\footnote{Usatges de Barcelona. Us. 21 (24--25): «PLACITVM MANDETVR tam magnatibus quam militibus, primum ad decem dies; deinde de octo in octo mandetur. Ad rusticos namque mandetur placitum in quarto die uel quinto». P.~68.}. Там же подчеркивается, что судить своих вассалов должен сеньор при своём дворе.

В 24 статье «Барселонских обычаев» сказано, что по обычным делам не должно быть более четырех судебных заседаний; первое отводится для изложения жалоб, второе --- для решения, выносимого судьями, избранными обеими сторонами, третье --- для его пересмотра, четвёртое --- для исполнения\footnote{Ibid. Us. 24 (28): «DE OMNIBVS NAMQVE comunibus causis non plus oportet quam quatuor esse placita: unum in quo sit directum firmatum per pliuios uel per pignora conuenienter sicut opus erit uel necesse, querimoniis ex utrisque partibus auditis; aliud namque in quo sint querimonie dicte et rationate, et iudicia data a iudicibus ex utrisque partibus electis; tercium quoque in quo sint a iudicibus querimonie et iudicia retracta, et, si opus erit uel necesse, iudicia meliorata: postea sint laudata et auctorizata et ad laudamentum iudicis illorum bene assecurata per pignora ut sint facta, et ibi debent crescere pignora ad laudamentum iudicis illorum; quartum namque in quo dominus placiti recuperet pignora, et, illo ea tenente, sint directa facta et iudicia completa, sicut erunt iudicata et ex utrisque partibus auctorizata». P.~70.}. Тем самым процесс предстает как поэтапная и формализованная процедура, а не как разовое судебное действие.

Нормы об обвинении подтверждают ту же тенденцию к процессуальной определенности. В статье B7 (90) «Барселонских обычаев» закреплено, что обвинение не должно приниматься в письменной форме и должно исходить от самого обвинителя: \emph{«никакое обвинение не должно приниматься по записи; пусть обвиняет собственным голосом»}\footnote{Ibid. Us. B7 (90): «Per scripturam nullius accusacio suscipiatur, sed propria uoce accuset, si legitima et condigna accusatoris persona fiat, presente uidelicet eo quem accusare desiderat, quia nullus absens aut accusari potest aut accusare». P.~170.}. Далее уточняется, что обвинитель должен быть надлежащим лицом, а обвиняемый --- присутствовать при обвинении, поскольку отсутствующий не может ни обвинять, ни быть обвиненным. Тем самым обвинение мыслится не как безличное сообщение о правонарушении, а как личное и публичное процессуальное действие, совершаемое в присутствии другой стороны и потому изначально включенное в пространство суда. В статье B5 (88) говорится: \emph{«Никто не должен одновременно выступать обвинителем, судьей и свидетелем, поскольку всякое судебное разбирательство требует обязательного присутствия четырех лиц…»}\footnote{Ibid. Us. B5 (88): «Nullus unquam presumat accusator simul esse et iudex et testis, quoniam in omni iudicio quatuor persone necesse est semper adesse…». P.~170.}. Суд, следовательно, требовал не только выдвижения обвинения, но и строгого разграничения процессуальных ролей. Это означает, что составители «Барселонских обычаев» уже достаточно отчетливо различали функции инициатора процесса, лица, выносящего решение, и лица, сообщающего сведения о факте. Иначе говоря, достоверность обвинения обеспечивалась не столько формальной письменной фиксацией, сколько публичностью заявления, личной ответственностью обвинителя и процессуальным разделением функций, исключавшим смешение интереса, свидетельства и судебной власти.

Наиболее разработанными в «Барселонских обычаях» являются именно положения о свидетелях, и это позволяет особенно ясно увидеть социальную природу средневекового доказательства. Статья B2 (85) предписывала произвести проверку свидетелей по репутационному критерию и отношению к сторонам процесса, а если иначе выяснить их пригодность невозможно,\emph{ «отделить друг от друга и расспросить поодиночке»}\footnote{Ibid. Us. B2 (85): «...separentur ab invicem et singulariter inquirantur...». P.~170.}. Здесь важен не только сам приём раздельного допроса, направленный против согласования показаний и действий по предварительному сговору, но и исходная установка на предварительную проверку свидетеля. Суд интересует не только содержание будущих показаний, но и то, кто именно говорит, из какой среды он происходит и насколько его слово заслуживает доверия. Поэтому в той же статье говорится, что свидетелей следует выбирать \emph{«из того же }\emph{пага}\emph{, а не из другого»}\footnote{Ibid.: «...separentur ab invicem et singulariter inquirantur...». P.~170.}: наиболее ценным признается свидетель, укорененный в той же локальной общине, где произошел спор или правонарушение. С точки зрения антропологии права это означает, что судебная истина мыслится как возникающая внутри социального контекста: знание о факте принадлежит тем, кто включен в пространство соседства, повседневного наблюдения и коллективной памяти\footnote{Подробнее о проблематике памяти в праве см. \emph{Клэнчи}\emph{ М.} Воспоминания о прошлом и старый добрый закон // Vox medii aevi / под ред. \emph{Горбун Г.С.}; пер. \emph{Горбун Г.С.}~2015. №~2--3 (13--14). С.~115.}.

\subsection{Число свидетелей и достаточность свидетельских показаний}

Эта же логика продолжается в статье B3 (86), где устанавливается, что больше доверия следует оказывать \emph{«более почтенным свидетелям, чем низшим»}\footnote{Usatges de Barcelona. Us. B2 (86): «...honestioribus magis quam vilioribus testibus fides pocius admittatur...». P.~170.}, а показание одного лица, даже если его персона кажется «пригодной», не должно приниматься само по себе\footnote{Ibid. Us. B3 (87): «...unius autem testimonium <…> nullatenus audiendum». P.~170.}. Это обстоятельство показывает нам, что уже составители «Барселонских обычаев» стремились ограничить репутационные формы судебного доказательства в пользу более рациональных и способствующих выяснению реальных обстоятельств дела. В статье B3 (87) закрепила в общем виде правило относительно числа свидетелей: \emph{«двух или трех пригодных свидетелей достаточно для доказательства всех дел»}\footnote{Ibid. Us. B3 (86): «...duo vel tres idonei testes ad omnia probanda negocia sufficiunt». P.~170.}, тогда как единичное свидетельство отвергается. Следовательно, свидетельство в барселонском праве строится на сочетании двух критериев: множественности и социальной квалификации свидетелей. Суд не удовлетворяется индивидуальным знанием факта, даже если оно исходит от лица высокого статуса; он требует, чтобы истина была подтверждена несколькими социально приемлемыми голосами. Тем самым правовой порядок опирается на коллективно удостоверенную достоверность, а не на изолированное заявление частного лица. Здесь обнаруживается прямая связь с репутационными формами доказательства: \emph{idoneus}\emph{ }\emph{testis} --- это не просто очевидец, а человек, уже признанный общиной достойным доверия.

Особенно показательна статья B6 (89), исключающая из числа свидетелей и обвинителей тех, кто \emph{«позавчера или третьего дня был врагом стороны, чтобы они не стремились вредить в гневе или мстить за обиду; следовательно, для суда необходима их неоскорбленная и неподозрительная расположенность»}\footnote{Ibid. Us. B6 (89): «Accusatores et testes esse non possunt qui ante esternum diem aut nudius tercius inimici fuerunt ne irati nocere cupiant, ne lesi se ulcisci uelint; inoffensus igitur accusatorum et testium affectus querendus est et non suspectus...». P.~170.}. Далее та же статья добавляет, что не считаются вполне пригодными и те свидетели, которым можно приказать стать свидетелями\footnote{Ibid.~: «...Idonei testes non uidentur esse quibus imperari potest ut testes fiant». P.~170.}. Эти положения позволяют увидеть, что в «Барселонских обычаях» свидетель определяется прежде всего через его социальные отношения. Для оценки его слова значимы не только статус и почтенность, но и отсутствие зависимости, вражды и внешнего принуждения. Иначе говоря, правовая ценность свидетельского показания производна от положения человека внутри сети социальных связей. Свидетель здесь не нейтральный носитель информации, а участник общины, чья речь становится доказательством лишь при условии, что она свободна от подозрительных аффектов, от мести, подчинения и локальной недобросовестности. Именно поэтому нормы о свидетелях в «Барселонских обычаях» тесно примыкают к репутационным формам установления истины: почтенность, местная укорененность, отсутствие вражды и независимость образуют ту социальную основу, без которой само сообщение о факте не является полноценным доказательством.

В «Обычаях Льейды», оказавших заметное влияние на последующую правовую традицию Новой Каталонии, еще не представлена столь разработанная регламентация розыскного судебного процесса, как в более поздних памятниках (например, «Обычаях Тортосы»), однако нормы о доказательствах и свидетельствовании уже являлись более строгими. Так, статья 131 устанавливает: \emph{«Погашение долга, удостоверенного публичным }\emph{документом, доказывается двумя достаточными свидетелями»}\footnote{Costumbres de Lleida. [131]: «…Probatur solucio crediti publici instrumenti per duos testes sufficientes...». P.~137.}. Норма исходила из принципа: кто утверждает факт, тот и должен его доказать. Неудача в доказательстве не давала стороне права переложить эту обязанность на противника в рамках того или иного спора.

Этой же логике соответствует и 133-я статья, содержащая отсылку к одной из норм «Вестготской правды»\footnote{LV. II.4.2: «…quia testes sine sacramento testimonium peribere non possunt» («…поскольку свидетели не могут свидетельствовать без клятвы»). P.~95; Подробнее о рецепции этой нормы в каталонском праве («Обычаи Тарреги», «Обычаи Балагера», «Обычаи Пералады») см. \emph{Font Rius J.M. }Costumbres de Tárrega // Anuario de Historia del Derecho Español. 1953. №~23. P.~441.}: \emph{«Свидетели не могут быть принуждены к даче показаний, кроме тех случаев, когда они сами обещали подтвердить обстоятельство, либо прямо названы в документах, либо специально выбраны свидетелями; тогда они обязаны либо свидетельствовать, либо поклясться в том, что им ничего не известно»}\footnote{Costumbres de Lleida. [133]: «Non coguntur testes nisi se promiserint probaturos, vel nisi in cartis scripti inveniantur, vel nisi eligantur ut sint testes, et tunc coguntur vel iurabunt se nichil scire». P.~137--138.}. Особенно важно и заключительное положение статьи: если свидетель не мог подтвердить обстоятельства, он был обязан не просто уклониться от участия, а официально поклясться, что ничего не знает. Это означает, что даже отказ от содержательного свидетельства должен был быть введен в процессуальную форму и получить юридическое выражение.

Наконец, 134 статья устанавливала, что для любого документа достаточно двух свидетелей, но в то же время любой документ может быть объявлен подозрительным, и такому возражению должно быть уделено внимание\footnote{Ibid. [133]: «In quolibet instrumento sufficiunt duo testes et in qualibet re vel causa. Suspectum potest quis dicere quodlibet instrumentum etiam sine inscripcione non abolitum, vel deletum, et est ei fides facienda». P.~138.}. Письменный акт в «Обычаях Льейды» еще не воспринимался как безусловно достаточное и само по себе решающее доказательство. Его доказательная сила зависела как от подтверждения со стороны свидетелей, так и от возможности процессуального оспаривания, а значит, окончательное значение документа определялось не только его письменной формой, но и его местом в общем ходе судебного разбирательства.

Таким образом, «Обычаи Льейды» фиксируют уже достаточно строгую иерархию видов доказательств, в которой свидетельские показания, документ и процессуальная обязанность сторон подчинены более строгим правилам, чем в ранних каталонских памятниках. Вместе с тем в «Обычаях Льейды» ни свидетельство, ни письменный акт еще не обладали полной автономией: их доказательная сила зависела от соблюдения процессуальной формы, от возможности оспаривания и от включенности в общий порядок судебного разбирательства. В этом отношении «Обычаи Льейды» занимают промежуточное положение между более ранней, менее расчлененной моделью процесса и позднейшими, более развитыми формами розыскного судопроизводства.

\subsection{Личное знание и разрешение противоречий между свидетелями}

В «Обычаях Тортосы» мы также можем зафиксировать эту тенденцию. Так, согласно статье 4.11.8: \emph{«Для доказательства в любом деле достаточно двух или более свидетелей, при условии, что они не пользуются дурной славой, либо публичного документа, либо завещания, либо иного последнего волеизъявления, удостоверенного двумя или более свидетелями»}\footnote{Costums de Tortosa. 4.11.8: «E totes coses a provar basten II testimonis o plus, ab que no sien de mala fama, o carta pública, o testament, o tota altra derrera volentat, en què aja II testimonis o plus». P.~208.}, что явно указывает на условное равенство письменного документа и свидетельских показаний в так называемой иерархии средств доказательства. Заметно, что ко 2-й половине XIII~в. надлежащим образом оформленный документ упрочил свои позиции в иерархии доказательств и стал весомой альтернативой такому способу доказательства, как клятва.

Одновременно «Обычаи Тортосы» запрещали свидетельство «с чужих слов» без личного подтверждения: \emph{«…всегда следует иметь в виду, }\emph{что свидетельство с чужих слов не имеет доказательной силы…»}\footnote{Ibid. 4.11.11: «…asò tota ora entés que testimoni d 'oyda ren no val…». P.~208.}. Такой запрет стал важным маркером смены культуры доказательства: внимание сместилось не только к количеству свидетелей, но и к качеству знания, получаемого в ходе их расспроса. Лишь непосредственные показания, подкрепленные торжественной клятвой, обрели полную юридическую силу.

Что касается процессуальных норм, зафиксированных в поздних сводах обычного права, таких как «Обычаи Орты» и «Обычаи Миравета», то их составители предпочитали практически полностью ориентироваться на вышеобозначенную модель. Рассмотрим несколько примеров из «Обычаев Орты»:

\begin{itemize}
\item норма о достаточности двоих свидетелей для полного доказательства (LX статья) \footnote{Costumbres de Orta. [LX]: «Duo testes legitimi sufficiant in qualibet re vel causa». P.~622.};
\item норма о возможности presumpcio veritatis (LXI статья)\footnote{\textbf{Presumpcio}\textbf{ }\textbf{veritatis}\textbf{ }--- презумпция истинности; неполное доказательство, которое делало для суда позицию той или иной стороны более вероятной, но требовало дополнительного подтверждения (см. \emph{Baldwin~J.W. }The Intellectual Preparation for the Canon of 1215 against Ordeals // Speculum. 1961. Vol.~4 (36). P.~616). Стоит также отметить, что «Обычаи Орты» предложили более чёткое разграничение полного и неполного доказательства в зависимости от предварительного заявления сторон: Costumbres de Orta. [LXI]: «Item, quod unus testis productus in presumpcionem veritatis sufficiat usque ad sumam centum solidos cum sacramento principalis, sit quod probans nisi dixerit tempore quo dicit se probatorum quod probabit in presumpcione veritatis videatur se strinxisse probare in plenam probacionem, et si unus testes hoc casum probaret, pro eo non obtineat nisi delato ei sacramento ubi modica quantitas pete- retur judicanti videretur aliud faciendum~». P.~622.}.
\end{itemize}
Однако ряд статей «Обычаев Орты» существенно отличался от аналогичных в «Обычаях Льейды»;\footnote{Ibid. [LXI]: «Mulieres discrete et bone fame possint facere testimonium in causis civilibus prout mares~». P.~622.}, это утверждение касается прежде всего статьи LXII о возможности женщин свидетельствовать в гражданских делах, в то время как «Обычаи Льейды» напрямую запрещали женское свидетельство практически полностью\footnote{Подробнее об этих нормах в контексте репутационного критерия см. §3.3 настоящего исследования.}.

В «Обычаях Миравета» порядок доказательства строился на нескольких уже знакомых нам принципах. Так, статья 107 гласит: \emph{«в любом документе достаточно двух свидетелей, равно как и по любому делу»}\footnote{Constituciones baiulie Mirabeti. 107: «In quocumque instrumento sufficiunt duo testes vel in qualibet re~». P.~28.}. Тем самым составители этого свода следовали той же модели, что и составители «Обычаев Льейды»: показания одного свидетеля сами по себе обычно не считались достаточными. Вместе с тем 106 статья показывала, что свидетельство понималось не только как право, но и как обязанность: \emph{«свидетель, который будет вызван для дачи показаний, должен быть принужден дать показание, с возмещением ему издержек и трудов»}\footnote{Ibid. 106: «Testis qui nominatus fuerit in testimonium [cogatur facere testimonium] satisfaciendo illi testi in suis laboribus~». P.~28.}. Это значит, что установление истины в суде уже зависело не только от частной инициативы сторон, но и от обязанности членов общины участвовать в правосудии. Несмотря на то что данная статья не была основана на нормах из «Обычаев Льейды» и носила оригинальный характер, в целом судебное доказательство и в этом кодексе сочетало в себе процессуальную определенность, репутацию участников и контролируемое использование свидетельских показаний.

\subsection{Вызов, принуждение и допрос свидетеля}

Социально-юридические условия допроса также были жестко определены в сводах местного обычного права. Так, согласно «Обычаям Тарреги» допрос могли проводить лишь специально уполномоченные чиновники, а именно вегер; кастеляны могли выступать инициаторами процесса розыскного типа лишь в исключительных случаях: \emph{«Кастеляны и их }\emph{байлы}\emph{ не вправе производить расследование по уголовному делу, если только преступление не было совершено в их присутствии. В таком случае они должны принять поручительство от лица, подавшего жалобу; если же после этого жалобщик захочет отказаться от обвинения, сеньор может по собственной должности продолжить расследование }\emph{преступления}»\footnote{Consuetudines ville Tarege. [3]: «Castlani nec illorum baiuli non faciant inquisicionem in aliqua causa criminali nisi in eorum presencia crimen fuerit comissum et tunc accipiant fideiussorem ab illo a quo querimonia facta fuerit et fideiussione receipta si querelator ab accusacione voluerit desistere, dominus ex officio suo possit de maleficio inquirere». P.~437.}. Эта норма показывает, что право стремилось ограничить возможность произвольного начала расследования со стороны местных властей и поставить допрос в рамки четко распределенной юрисдикции. Вместе с тем она фиксирует важный сдвиг от частно-обвинительной логики к служебному преследованию: даже если жалобщик отказывался от обвинения, само дело не прекращалось автоматически, поскольку сеньор сохранял возможность продолжить его самостоятельно. Тем самым допрос и расследование в Тарреге уже мыслились не только как средство защиты частного интереса потерпевшего, но и как инструмент поддержания публичного порядка.

Исследование регламента допроса в каталонских источниках середины --- 2-й пол. XIII~в. показывает переход к инквизиционной процедуре сбора доказательств: ключевую роль начинают играть сбор улик, письменные документы и опрос свидетелей. Так, например, «Кутюмы Перпиньяна» показывают, что судопроизводство строилось на сочетании нескольких принципов: сравнительно высокой доказательной силы письменного акта, довольно четкой фиксации процессуальной позиции сторон и коллегиального контроля над судебными решениями. Так, в статье LXVII говорилось, что если в ходе процесса по промежуточному решению (\emph{interlocutoria}) испрашивается его «улучшение» или пересмотр (\emph{melioramentum}), то суд может в тот же или в другой день подтвердить либо отменить это решение, но лишь после нового совещания с \emph{«другим пригодным советом добрых людей, а не того, кто вынес это решение»}\footnote{Les Coutumes de Perpignan. [LXVII]~: «…curia <…> confirmare vel infirmare interlocutoriam <…> dum per alium ydoneum faciat consilio proborum hominum, non illius qui illam protulit». P.~585--586.}. Иначе говоря, промежуточное судебное постановление не считалось окончательно закрытым для пересмотра, однако такой пересмотр должен был осуществляться не единолично и не тем же составом советников, а с привлечением иных компетентных лиц. Это позволяет увидеть в судебной процедуре стремление не только к оперативности, но и к формальной беспристрастности.

Согласно королевской привилегии «Знатные мужи Барселоны постановили…», розыскной процесс уже являлся довольно строго ограниченной процедурой. Прежде всего, в тексте прямо установлено, что вегеры и другие королевские должностные лица не должны проводить против граждан Барселоны ни общего, ни специального расследования, кроме случаев уголовного преступления. Более того, такое расследование должно вестись, согласно 100 статье, \emph{«с одним знатоком права и двумя добрыми людьми»}\footnote{Recognoverunt proceres. [C]: «…cum uni iurisperito et duobus probis hominibus…». P.~610.}, то есть с участием не только профессионального юриста, но и представителей городской общины. Эта же линия прослеживается и в следующей норме: вегеры и байлы не могли поручать судебные дела или расследования лицам, «\emph{которые не являются юристами»}\footnote{Ibid. [CIV]: «...quod vicarii vel baiuli non possint comitere iudicia vel inquisitiones aliquibus qui non sint iurisperiti...». P.~611.}. Следовательно, барселонская модель розыскного процесса сочетала два принципа: с одной стороны, профессионализацию процедуры, с другой --- ее включенность в городской правовой порядок через участие знати, байлов и вегеров.

Мы видим, что розыскной тип процесса применялся исключительно по отношению к уголовным преступлениям и был поставлен под контроль юридически компетентных лиц и вместе с тем ограничивался рядом процедурных гарантий, прежде всего в вопросах расследования против горожан и принуждения к свидетельству.

\section{Письменный акт как судебное доказательство}

Письменные акты занимали в каталонском судебном разбирательстве XII--XIII~вв. неодинаковое положение. Одни из них составлялись до возникновения конфликта и впоследствии использовались для подтверждения приобретения имущества, существования обязательства или содержания последней воли; другие возникали непосредственно в ходе тяжбы либо закрепляли её результат. При этом доказательная сила документа определялась не только его содержанием и письменной формой, но также способом удостоверения, участием свидетелей, происхождением из признанной канцелярии и возможностью проверить обстоятельства его составления.

Задача настоящего параграфа состоит в том, чтобы установить основные функции письменных актов при разрешении тяжб и определить условия, при которых они могли подтверждать право или обстоятельства дела. Для этого сначала будут рассмотрены виды документов и их место в развитии конфликта, затем --- соотношение письменного акта со свидетельским показанием и значение публичного удостоверения. Наконец, необходимо проследить порядок предъявления и оспаривания документов, а также способы сохранения подтверждённого ими права при утрате первоначальной грамоты или возобновлении уже разрешённого спора.

\subsection{Виды письменных актов и их функции в тяжбе}

Письменные акты, упоминаемые в каталонских правовых памятниках XII--XIII~вв., не образовывали единой категории. Они различались по содержанию, способу удостоверения и обстоятельствам составления, а главное --- по той функции, которую могли выполнять в рамках потенциальной тяжбы. Одни документы появлялись до конфликта и позднее предъявлялись как основание имущественного или обязательственного требования. Другие создавались уже в ходе разбирательства: фиксировали жалобу, показания сторон и свидетелей либо принятое решение. Наконец, отдельную группу составляли грамоты, которыми закреплялись примирение, признание права противной стороны или окончательный отказ от притязаний. Поэтому письменный акт следует рассматривать с учётом его назначения и обстоятельств составления.

В «Барселонских обычаях» письменные акты выступали основанием ранее приобретённого права. Так, согласно статье 54(57), рыцарь, требовавший передачи феода, которым он в тот момент не владел, должен был подтвердить, что получил его от сеньора, «свидетелями или грамотами»\footnote{Usatges de Barcelona. Us. 54(57): «Illos autem quos non tenuerint et clamauerint, aut probent per testes uel per scripturas eos a senioribus eorum adquisisse, aut dimittant eos». P.~90--91.}. В данном случае источник не определял вид предъявляемого документа, обозначая его общим словом «грамоты» (лат. --- \emph{scripturae}; кат. --- \emph{escriptures}). Их функция, однако, очерчена достаточно ясно: письменный акт подтверждал совершённую ранее передачу феода и тем самым служил основанием заявленного требования.

Более определённо разновидности письменных актов названы в статье C3 (148) того же памятника. В тяжбе могла быть предъявлена \emph{«завещательная грамота или надлежащим образом скреплённая грамота относительно предмета спора»}\footnote{Ibid. C3 (148): «...testamentum vel carta firmata de aliqua contencione...». P.~171.}. Здесь письменный акт выполнял уже две возможные функции. Завещание подтверждало распоряжение имуществом на случай смерти, тогда как содержание «скреплённой грамоты» могло относиться и к иному правовому действию, ставшему основанием спора. Противная сторона, в свою очередь, могла противопоставить предъявленному акту другие надлежащим образом удостоверенные письменные документы (лат. --- \emph{firmae}\emph{ }\emph{scripturae}). Таким образом, «Барселонские обычаи» предусматривали использование грамот для подтверждения различных имущественных прав, в том числе возникших из завещательного распоряжения; такие документы могли служить основанием требования или возражения.

В «Обычаях Льейды» письменные акты упоминались преимущественно в связи с обязательствами и распоряжением имуществом. Наиболее ясно выделялся публичный акт о долге (лат. --- \emph{publicum}\emph{ }\emph{crediti}\emph{ }\emph{instrumentum}), подтверждавший возникновение и содержание обязательства. Наряду с ним текст предполагает существование письменного подтверждения произведённой выплаты, а также грамот, оформлявших передачу земли в держание. Письменный акт, таким образом, мог удостоверять не только возникновение права или обязанности, но и их последующее изменение или прекращение. Особенно показательно, что источник разграничивал сам долг, закреплённый в публичном акте, и факт его погашения: \emph{«Погашение долга, удостоверенного публичным актом, доказывается двумя достаточными свидетелями»}\footnote{Costums de Lleida. [131] // Villanueva J. Viage literario á las iglesias de España. T. XVI: Viage á Lérida. Madrid: Imprenta de la Real Academia de la Historia, 1851. Apéndice: «Probatur solucio crediti publici instrumenti per duos testes sufficientes…». P.~187.}. В этой норме акт подтверждал существование обязательства, тогда как его исполнение представляло собой отдельное обстоятельство, требовавшее самостоятельного удостоверения.

В памятниках, созданных на основе «Обычаев Льейды», указания на подобные документы получали более конкретное содержание. Согласно 11 статье «Обычаев Тарреги», при продаже, отчуждении или иной передаче имущества, которым держатель владел на праве эмфитевзиса, он должен был выплатить сеньору пятидесятую часть цены, а сеньор --- подтвердить составленную грамоту\footnote{Consuetudines ville Tarege. [11]: «Emphiteota si vendiderit, alienaverit vel ad aliquem transtulerit rem quam pro domino teneat, donet quinquagesimam partem precii domino pro quo rem tenuerit et firmet cartam dominus nisi amplius vel paucis ex conventione inhita inter partes statutum fuerit, exceptis tamen hiis qui pro nobis tenuerint». P.~439.}. После уплаты сеньору причитавшейся ему доли он должен был подтвердить соответствующий акт. Такая грамота закрепляла одновременно отчуждение имущества и согласие лица, за которым сохранялись сеньориальные права. Источник не описывал её непосредственного предъявления в тяжбе, однако соблюдение письменной формы и подтверждение сеньором создавали основание для последующей защиты приобретённого права.

В «Обычаях Орты» письменные акты выступали прежде всего основанием коллективных прав и имущественных обязанностей жителей. Несколько статей отсылали к ранее пожалованным грамотам, закреплявшим права общины на дороги, улицы, площади, пастбища, воды и проведение рынка. Такие документы не только свидетельствовали о предшествовавшем пожаловании, но и определяли объём прав, признаваемых за жителями в конце XIII~в. Так, право пользоваться дорогами, улицами и площадями сохранялось \emph{«в соответствии с актами, пожалованными им братьями Храма»}\footnote{Costumbres de Orta. [XIII]: «Item, quod dicti homines habeant vias, vicos et plateas in Orta et aliis locis termini de Orta, et inde possint ire, redire et manere libere et spaciose, prout in instrumentis concessis eis a fratribus Templi continetur. Item, quod mercatum ville de Orta non mutetur de loco et die, prout in instrumentis factis super dicto mercato eisdem hominibus continetur». P.~619}. Сходным образом место и день проведения городского торга определялись ранее составленными по этому поводу грамотами.

Отдельно «Обычаи Орты» упоминали письменный акт, определявший обязанности должника перед орденом тамплиеров. Если должник или поручитель не производил выплату в установленный срок, комендатор либо байл должен был взыскать долг и причитавшуюся ордену треть \emph{«в соответствии с тем, что содержится в акте, составленном для жителей братьями Храма»}\footnote{Ibid. [VII]: «…comendator vel bajulus compellat solvere debitorem illud debitum et tercium Templi, prout in instrumento a fratribus Templi super hoc eis facto continetur». P.~618.}. В этом случае документ устанавливал размер и порядок исполнения обязательства и служил основанием действий сеньориальной администрации.

Более поздние «Обычаи Миравета» не содержали развёрнутого перечня разновидностей документов, но сохраняли общее обозначение письменного акта (лат. --- \emph{instrumentum})\footnote{Constituciones baiulie Mirabeti. 107: «In quocumque instrumento sufficiunt duo testes vel in qualibet re». P.~28.}. Их значение для настоящего сюжета состоит прежде всего в том, что они воспроизводили представление о письменном документе как об акте, который должен был быть удостоверен свидетелями. Поскольку этот памятник относится уже к началу XIV~в., его следует использовать не как основной источник для рассматриваемого периода, а как позднее свидетельство устойчивости норм, восходивших к «Обычаям Льейды».

В «Обычаях Тортосы» перечень письменных актов, способных иметь значение при разрешении тяжбы, получил более отчётливое выражение. В статье 4.11.8 наряду со свидетельскими показаниями названы публичная грамота, завещание и иное последнее волеизъявление: \emph{«Для доказательства любых обстоятельств достаточно двух или более свидетелей, при условии, что они не пользуются дурной славой, либо публичной грамоты, либо завещания, либо иного последнего волеизъявления, при котором присутствовали два или более свидетеля»}\footnote{Costums de Tortosa. 4.11.8: «E totes coses a provar basten II testimonis o plus, ab que no sien de mala fama, o carta pública, o testament, o tota altra derrera volentat, en què aja II testimonis o plus». P.~208.}.

Перечисление позволяет различить три разновидности письменного удостоверения. Публичная грамота (кат. --- \emph{carta}\emph{ }\emph{pública}) могла относиться к имущественному или обязательственному действию и обладала особым значением благодаря порядку своего составления. Завещание (кат. --- \emph{testament}) представляло собой отдельный оформленный акт. Наконец, иное последнее волеизъявление (кат. --- \emph{tota}\emph{ }\emph{altra}\emph{ }\emph{derrera}\emph{ }\emph{volentat}) могло первоначально существовать в устной форме, но приобретало устойчивое правовое значение после его записи и подтверждения свидетелями. Таким образом, «Обычаи Тортосы» учитывали не только готовые письменные документы, но и случаи последующего письменного удостоверения ранее произнесённой последней воли.

В «Кутюмах Перпиньяна» особое положение занимал публичный акт, составленный в городской канцелярии. Согласно статье LXII, лицо, требовавшее выплаты денег на основании такого документа, не обязано было подробно излагать основание своего требования: достаточно было предъявить \emph{«публичный акт, составленный в писчей канцелярии Перпиньяна»}\footnote{Les Coutumes de Perpignan. [LXII]: «…cum instrumento publico facto in scribania Perpinyani». P.~32.}. Письменный акт здесь влиял на форму заявления требования, освобождая истца от необходимости подробно воспроизводить обстоятельства его возникновения.

Дополнить перечень категорий документов, связанных с имущественными отношениями, позволяет королевская привилегия «Знатные мужи Барселоны постановили…» (Recognoverunt proceres). В статье LVI приводилась формула брачной грамоты, в которой одной суммой закреплялись приданое женщины и предоставляемое мужем приращение к нему: \emph{«Когда грамота составляется в следующих словах: “Я, такой-то, даю тебе, такой-то, моей жене, такую-то сумму в качестве твоего приданого и приращения к нему”, следует понимать, что одна треть этой суммы является дарением по случаю брака»}\footnote{Recognoverunt proceres. [LVI]: «Item, quando instrumentum dotalium fit in hec verba: “Ego, talis, dono tibi tali uxori mee, tantum pro dote tua et sponsalicio”, quod intelligatur tercia pars ipsius quantitatis esse donatio propter nuptias$\rightarrow$Encara con la carta de l’espoali se ffa en estes paraules: “Jo aytal do a tu aytal, muller mia, aytant per ton exovar e per ton espoali”. Que la terça part d’aquella cantitat es entesa que es creix». P.~601.}. Таким образом, статья не только приводила образец составления брачной грамоты, но и устанавливала порядок толкования указанной в ней совокупной суммы: две трети приходились на приданое, а одна треть представляла собой дар мужа по случаю брака. Такая грамота могла подтверждать размер имущественных требований женщины к имуществу мужа, хотя сама статья регулировала прежде всего содержание и толкование акта, а не порядок его предъявления в суде.

Нормативные памятники, следовательно, позволяют выделить несколько основных групп письменных актов, которые могли приобрести значение при возникновении тяжбы. К ним относились:

\begin{itemize}
\item грамоты, подтверждавшие передачу имущества или феода;
\item документы о возникновении, исполнении или прекращении обязательств;
\item завещания и записи иной последней воли;
\item брачные грамоты, закреплявшие приданое и брачный дар;
\item публичные документы городских канцелярий;
\item наконец, грамоты о пожаловании, подтверждавшие коллективные права общины. Однако сам факт упоминания документа в нормативном тексте ещё не означает, что источник описывает его непосредственное предъявление суду. В одних случаях прямо зафиксировано использование акта в тяжбе, в других письменная форма лишь создавала основание права, которое потенциально могло нуждаться в защите.
\end{itemize}
Актовый материал позволяет точнее установить функции этих документов. Особенно ясно письменный акт как основание требования представлен в грамоте 1017~г. о споре между монастырём Сан-Кугат и вегером Сунифредом де Руби. Во время разбирательства аббат Гитард предъявил два уже существовавших документа: грамоту о продаже спорных аллодов его предшественнику и грамоту об их дарении монастырю. В тексте отмечено, что аббат \emph{«предъявил составленные грамоты: одну о продаже <…>, а другую о дарении»}\footnote{Justícia i resolució… 2018. Doc. 174: «…protulit iam dictus abbas scripturas confectas; unam venditionis <…> et alteram donationis». P.~336.}, благодаря которым суд установил принадлежность спорного имущества монастырю. Грамоты о продаже и дарении были составлены до возникновения тяжбы, но в суде получили новую функцию --- стали основанием заявленного права.

Сходное сочетание нескольких письменных оснований обнаруживалось и в споре 1156/1157~г. о замке Гаварра. Виконт Жерал де Кабрера признавал, что его дед передал замок церкви посредством дарственной грамоты, а позднее вновь распорядился им в пользу той же церкви в своём завещании: замок был подарен \emph{«посредством дарственной грамоты»}, а затем оставлен церкви по завещанию\footnote{Justícia i resolució … 2024. Doc. 328: «…per scripturam donationis et postea illud eidem ecclesie in testamento suo dimisit…». P.~375.}. Здесь дарственная грамота и завещание образовывали последовательную цепь письменных оснований, подтверждавших право церкви.

Наряду с документами, составленными до начала конфликта, в актовом корпусе представлены тексты, возникавшие непосредственно в ходе разбирательства. Одним из наиболее характерных видов были перечни жалоб и притязаний (лат. --- \emph{querimoniae}). Составители корпуса XII~в. определили их как неофициальные памятные записи, обычно не имевшие даты и подписей. Они конкретизировали обвинения, предъявляемые перед судьями или арбитрами, и могли служить основой последующих переговоров с противной стороной. Таким образом, \emph{querimonia} не подтверждала ранее приобретённое право сама по себе, а письменно оформляла позицию стороны и вводила отдельные претензии в пространство судебного или арбитражного разбирательства.

В ходе тяжбы также могли составляться записи свидетельских показаний, судебные решения, акты осмотра спорных владений и обязательства сторон подчиниться будущему решению. Эти документы не существовали до конфликта и не являлись первоначальным основанием требования. Их функция состояла в письменном закреплении отдельных процессуальных действий и сведений, полученных в ходе разбирательства.

Наконец, особую группу составляли акты, которыми оформлялось завершение конфликта. В корпусах IX--XII~вв. они обозначались как признание, отказ от права, окончательное прекращение притязаний или примирение (лат. --- \emph{recognitio}, \emph{evacuatio}, \emph{definitio}, \emph{pacificatio}). Составители корпуса отмечают, что подобные грамоты нередко представляли собой документы об исполнении решения: сторона отказывалась от спорного права либо признавала его за противником, а письменный акт закреплял достигнутый результат. С первой трети XI~в. к более ранним формулам отказа и признания всё чаще присоединялось указание на примирение (лат. --- \emph{pacificatio}).

Показательна грамота от 28 февраля 1157~г., в которой Беренгер де Англес отказался в пользу Жиронской епископальной кафедры, представленной епископом Беренгером де Льерсом и его преемниками, от поборов и иных притязаний в отношении церквей Святого Христофора в Лес-Планес, Святой Марии в Лез-Энсиес, Святого Ацискла в Кольторте и Святого Викентия в Камосе. Он заявлял, что \emph{«окончательно отказывается, освобождает и всеми способами оставляет»}\footnote{Justícia i resolució… 2024. Doc. 327: «…diffinio et evacuo et omnibus modis derelinquo...». P.~374.} прежние притязания, а сам акт был назван \emph{«грамотой об окончательном отказе»}\footnote{Ibid.~Doc. 327:«...scriptura diffinitionis...». P.~375.}.

Наряду с документами, составленными до возникновения конфликта, в актовом корпусе представлены тексты, фиксировавшие позицию стороны уже после начала спора --- на подготовительной стадии или непосредственно в ходе его разрешения. К ним относились перечни жалоб и притязаний (лат. --- \emph{querimoniae}). Составители издания корпуса документов XII~в. характеризуют их как неофициальные памятные записи, обычно не имевшие даты и подписей, которые конкретизировали обвинения перед судом, чаще всего арбитражным, и могли использоваться при переговорах с противной стороной.

Так, в перечне жалоб Гитарда Исарна, сеньора Кабоэта, датированном приблизительно ноябрём 1105~г., последовательно перечислялись нарушения, в которых он обвинял своего вассала Гильема Арнау, его сына и Миро Арнау: распоряжение переданным им имуществом вопреки установленным условиям, отказ восстановить право сеньора, неисполнение феодальных обязательств, присвоение имущества и оскорбления его жены. Запись завершалась обращением к байлам и верным помочь его жене и сыну вернуть принадлежавшие им права \emph{«войной или правосудием»}\footnote{Ibid. Doc. 17: «…per pled et per gera tro lor dret los en sia exid…». P.~23.}.

Ещё нагляднее связь такого перечня с последующим разбирательством прослеживается в документах 527--528. Между 1164 и 1173~гг. граф Арнау Миро I перечислил свои притязания к Рамону д’Эрилю: нарушение перемирия, убийство людей графа, захват овец, неисполнение служб, невыплату долгов, захват Валь-д’Арана, причинение ущерба жителям долины Бенаск и препятствия деятельности рынка в Кастельо-д’Энкус. Перечень завершался списком заложников, предоставленных в обеспечение явки Рамона д’Эриля и его подчинения судебному решению\footnote{Ibid. Doc. 527. P.~643--644.}. Следующий акт открывался словами: \emph{«Это судебное решение по жалобам между графом }\emph{Пальярса}\emph{ и Рамоном }\emph{д’Эрилем}\emph{»}\footnote{Ibid. Doc. 528: «Hoc est iudicium de querimoniis que sunt inter comitem Paliarensem et Raimundum de Eril». P.~645.}, после чего по отдельным обвинениям приводились решения суда. Перечень претензий не удостоверял ранее приобретённое право сам по себе, но письменно формулировал притязания стороны и могла становиться непосредственной основой их судебного рассмотрения.

Таким образом, письменные акты, связанные с разрешением тяжб, различались прежде всего по своим функциям. Одни удостоверяли приобретение имущества, возникновение обязательства, завещательное распоряжение или привилегию общины и могли позднее быть предъявлены как основание требования. Другие письменно оформляли жалобу, позицию стороны либо сведения, полученные в ходе разбирательства. Третьи закрепляли принятое решение, примирение или отказ от дальнейших притязаний. Между этими группами не существовало чёткой границы: документ, первоначально оформлявший продажу, дарение или последнюю волю, позднее мог стать доказательством в тяжбе, тогда как грамота, составленная по итогам конфликта, в будущем служила основанием для отклонения возобновлённых требований. Поэтому значение письменного акта определялось не только его формой и содержанием, но и тем, на каком этапе конфликта он был составлен и какое действие должен был произвести в суде.

\subsection{Письменный акт и свидетельское показание}

Распространение письменных актов не привело к отделению документа от свидетельского показания. Напротив, каталонские правовые памятники показывают, что письменная запись и слово свидетеля могли выполнять по отношению друг к другу несколько различных функций. В одних случаях они выступали как альтернативные основания подтверждения права; в других свидетели участвовали в удостоверении самого акта и впоследствии должны были подтвердить зафиксированное в нём действие; наконец, свидетельское показание позволяло установить обстоятельства, которые возникли уже после составления документа или не могли быть удостоверены самим письменным актом. Поэтому соотношение документа и свидетельского слова нельзя свести к постепенному вытеснению одного другим.

Как было показано выше, уже «Барселонские обычаи» допускали подтверждение права на феод \emph{«свидетелями или грамотами»}\footnote{Usatges de Barcelona. Us. 54(57): «…per testes uel per scripturas…». P.~90--91.}, а в «Обычаях Тортосы» два или более свидетеля были названы наряду с публичной грамотой среди достаточных оснований для установления обстоятельств дела\footnote{Costums de Tortosa. 4.11.8: «E totes coses a provar basten II testimonis o plus, ab que no sien de mala fama, o carta pública, o testament, o tota altra derrera volentat, en què aja II testimonis o plus». P.~208.}. В подобных нормах документ и показания свидетелей представлены как разные способы подтверждения одного и того же права или события. Однако их присутствие в одних и тех же статьях ещё не означало, что письменный акт существовал независимо от свидетельского удостоверения. Во многих случаях именно присутствие свидетелей при совершении правового действия придавало записи значение, позволявшее позднее обратиться к ней в тяжбе.

Особенно ясно эта связь выражена в «Обычаях Льейды». Согласно статье 133, свидетели обычно не могли быть принуждены к даче показаний, кроме случаев, когда они заранее обещали подтвердить определённое обстоятельство, были названы в грамотах либо специально выбраны для свидетельствования: \emph{«Свидетелей не принуждают к даче показаний, если только они не обещали подтвердить обстоятельство, }\textbf{\emph{не были записаны в грамотах}}\emph{ или не были избраны свидетелями}\emph{. Т}\emph{огда они обязаны либо }\emph{дать показания, либо поклясться, что ничего не знают»}\footnote{Costumbres de Lleida. [133]: «Non coguntur testes nisi se promiserint probaturos, vel nisi in cartis scripti inveniantur, vel nisi eligantur ut sint testes, et tunc coguntur vel iurabunt se nichil scire». P.~188.}. Указание имени человека в грамоте имело последствия, выходившие за пределы момента её составления. Запись сохраняла сведения о лицах, в присутствии которых было совершено действие, а в случае тяжбы позволяла потребовать от них его подтверждения. При составлении акта человек выступал свидетелем продажи, дарения, признания долга или иного действия. В суде же он давал показания о том, что видел или слышал при его совершении. Письменный акт соединял эти два момента, превращая имя свидетеля в указание на лицо, к знанию которого можно было обратиться после возникновения спора.

Эта норма не означала, что свидетель, имя которого было указано в тексте грамоты, был обязан безусловно подтвердить её содержание. Если он не располагал необходимыми сведениями, то должен был поклясться, что ничего не знает. Тем самым письменная запись устанавливала обязанность явиться и ответить за подпись, но не заменяла личного знания свидетеля и не предопределяла содержание его показаний. Документ указывал на предполагаемого носителя сведений, тогда как их наличие должно было быть подтверждено самим вызванным лицом.

Следующая статья «Обычаев Льейды» непосредственно связывала действительность письменного акта с определённым числом участвовавших в его удостоверении лиц\footnote{Ibid. [134]: «In quolibet instrumento sufficiunt duo testes et in qualibet re vel causa». P.~188.}. В данном случае два свидетеля не противопоставлялись документу как другое средство судебного доказательства, а входили в порядок его составления и удостоверения. Правовое значение записи основывалось не только на сохранённом тексте, но и на возможности связать его с лицами, присутствовавшими при совершении зафиксированного действия.

Сходное установление воспроизведено в более поздних «Обычаях Миравета». Письменный акт в них не противопоставлялся свидетельскому слову: присутствие двух свидетелей входило в число условий, позволявших использовать запись для подтверждения права\footnote{Constituciones baiulie Mirabeti. 107: «In quocumque instrumento sufficiunt duo testes vel in qualibet re». P.~28.}.

Иной вариант взаимодействия документа и свидетельского показания возникал тогда, когда они относились к различным этапам одного правового отношения. Публичный акт мог удостоверять возникновение обязательства, но не содержал сведений о его позднейшем исполнении. Поэтому «Обычаи Льейды» устанавливали, что погашение долга, закреплённого публичным актом, подтверждается двумя достаточными свидетелями\footnote{Costumbres de Lleida. [131]: «Probatur solucio crediti publici instrumenti per duos testes sufficientes». P.~187.}. Документ сохранял сведения о существовании долга, тогда как свидетели сообщали о последующем событии --- передаче денег или ином исполнении обязательства.

Это различие важно для понимания функции письменного акта. Он не создавал неподвижную картину отношений сторон и не мог сам по себе показать, сохранялось ли записанное обязательство к моменту тяжбы. Если долговая грамота оставалась у кредитора после выплаты, должник должен был противопоставить ей сведения о состоявшемся исполнении. Свидетельское показание в этом случае не удостоверяло документ и не опровергало точность его первоначального содержания. Оно подтверждало новое обстоятельство, вследствие которого зафиксированное в акте требование более не подлежало исполнению.

Вместе с тем не всякое свидетельское сообщение могло быть противопоставлено письменному акту с одинаковым результатом. Статья 129 «Обычаев Льейды» предусматривала, что если кредитор на основании грамоты требовал передать заложенную вещь, а её владелец выдвигал возражение, способное устранить требование, но не подтверждал его, кредитор приносил клятву; при этом «\emph{презумпция, основанная на показании одного свидетеля, против письменного акта не допускается»}\footnote{Ibid. [129]: «...unius testis presumpcio contra instrumentum non admittitur». P.~187.}.

Следовательно, документ мог обладать большей силой, чем сообщение одного свидетеля. Однако и эта норма не устанавливала безусловного превосходства письменного акта: она исключала лишь презумпцию, основанную на единственном свидетельском показании, тогда как другие способы подтверждения возражения сохранялись. Более того, кредитор должен был подкрепить своё требование клятвой, если противная сторона заявляла об обстоятельстве, препятствовавшем его удовлетворению, но не могла это обстоятельство подтвердить. Письменный акт усиливал позицию кредитора, однако не устранял необходимость дополнительных процессуальных действий: допроса должника или привлечения свидетелей.

Особенно отчётливо взаимозависимость записи и свидетельского знания проявлялась при утрате документа. В таких случаях показания использовались не для подтверждения отдельного правового действия, а для восстановления содержания самой грамоты. Показателен акт, составленный в связи с утратой графиней Герсендой грамоты о продаже аллодов. Пресвитер Сесульд сообщил, что неоднократно читал документ, видел знаки продавца и других лиц, скрепивших его, и знал всю его удостоверительную часть. Другие свидетели заявили, что неоднократно слышали чтение грамоты и также знали её содержание: «\emph{Я, пресвитер }\emph{Сесульд}\emph{, неоднократно читал и перечитывал её <…> и достоверно знал всё содержание удостоверительной части этой грамоты, поскольку видел её собственными глазами. Мы же, другие упомянутые свидетели, <…> слышали, как её читали и перечитывали, и достоверно знали всё содержание удостоверительной части этой грамоты»}\footnote{Justícia i resolució… 2018. Doc. 68: «Et ego Sesuldus presbiter eam legi et relegi plures vices <…>. Et omnem testum firmitatis ipsius scripture firmiter novi oculis meis videndo. Et nos alii predicti testes <…> eam audivimus legentem et relegentem et omnem testum firmitatis ipsius scripture firmiter novimus». P.~152.}.

В этом случае предметом свидетельского показания было не только совершение продажи. Свидетели воспроизводили сведения о самом документе: его содержании, составителе, знаках продавца и лиц, участвовавших в удостоверении. Одни знали грамоту посредством непосредственного чтения, другие --- благодаря неоднократному публичному оглашению. Графиня, в свою очередь, поклялась, что документ действительно был утрачен и что она не скрывает его обманным путём. Тем самым восстановление письменного акта опиралось на соединение трёх элементов: памяти свидетелей о тексте, знания ими порядка его удостоверения и клятвы лица, заявившего об утрате.

Подобная практика соответствовала положению «Вестготской правды», допускавшему восстановление утраченной грамоты по показанию ещё живого свидетеля, подписавшего её, а в случае его смерти --- по показаниям других правоспособных и осведомлённых свидетелей, которые видели документ и в полной мере знали его текст или удостоверительную часть\footnote{LV. VII.5.2~: «Si testis, qui in eadem scriptura suscripsit, adhuc supprestis existit, per ipsum poterit coram iudice omnis ordo scripture perdite reparari. Quod si testem ipsum, qui in eadem scriptura suscriptor accessit, mortuum esse contigerit, tunc si legitimi et cognitiores repperti fuerint alii testes, qui eandem scripturam se dicant vidisse et omnem textum vel firmitatem eiusdem scripture plenissime nosse, similiter publica iudicum investigatione per eorum testimonium ille, qui scripturam perdidit, poterit suam reparare et percipere veritatem». P.~305. См. также: \emph{Сафронова Д.П.} Фальшивые грамоты в леонской правовой практике X--XI веков // История: переводить, понимать, оценивать: к юбилею \emph{М.А.~Юсима}. М.: Институт всеобщей истории РАН, 2021. С.~126.}. Приведённая грамота показывает, каким образом это правило действовало на практике: свидетели должны были не просто заявить, что грамота когда-то существовала, а воспроизвести её содержание и указать признаки, свидетельствовавшие о надлежащем составлении.

Здесь письменное и устное не являлись двумя независимыми носителями одних и тех же сведений. Свидетельская память обеспечивала продолжение действия документа после исчезновения его материального носителя. Вместе с тем она была организована вокруг письменного текста: свидетели вспоминали не только совершённое действие, но и саму грамоту, её формуляр, знаки удостоверения и неоднократное чтение. Письменный акт, таким образом, формировал содержание последующего свидетельского сообщения, тогда как свидетельское сообщение позволяло сохранить правовое значение утраченной записи.

Рассмотренные нормы и казусы позволяют выделить несколько вариантов взаимодействия документа и свидетельского показания. Они могли служить альтернативными основаниями подтверждения одного права; свидетели могли участвовать в удостоверении акта и позднее подтверждать совершённое при них действие; их показания могли относиться к обстоятельствам, возникшим после составления документа; наконец, через свидетельскую память восстанавливались содержание и порядок удостоверения утраченной грамоты. При этом письменный акт не просто сосуществовал со свидетельским словом. Он определял круг лиц, к которым можно было обратиться в случае спора, разграничивал первоначальное право и его последующие изменения и позволял сохранять сведения на протяжении времени. Свидетельское показание, в свою очередь, связывало запись с совершённым действием и восполняло те обстоятельства, которые не могли быть установлены по одному только тексту.

Следовательно, распространение письменных актов не устранило персонального характера удостоверения права. Запись делала сведения более устойчивыми, но её действие во многом поддерживалось лицами, присутствовавшими при составлении документа, названными в нём или знавшими его содержание. Изменилось не столько значение свидетельского слова само по себе, сколько его место по отношению к письменному акту: свидетель всё чаще подтверждал не право вообще, а составление документа, последующее изменение зафиксированного в нём отношения или содержание утраченной записи.

\subsection{Публичное удостоверение и условия признания письменного акта}

Связь письменного акта со свидетелями сама по себе не объясняет, почему отдельные документы получали в суде особое значение. В городских правовых памятниках XIII~в. наряду с присутствием свидетелей всё более важным становилось происхождение документа: кем он был составлен, в чьём присутствии, от имени какого учреждения и сохранялась ли соответствующая запись в городской канцелярии или судебном реестре. Публичное удостоверение не следует понимать как простое придание тексту известности. Оно предполагало включение письменного акта в установленный порядок деятельности писца, нотариуса, суда или городской канцелярии, благодаря чему достоверность документа связывалась уже не только с памятью частных лиц, но и с признанным правовым положением составившего его должностного лица.

Наиболее последовательно различие между публичной грамотой и иными письменными актами проведено в «Обычаях Тортосы». В рубрике о доказательствах устанавливалось, что все факты и договоры могли быть подтверждены признанием противной стороны, свидетелями или публичными грамотами:\emph{ «В городе Тортосе все факты и все договоры доказывают тремя способами: признанием противной стороны, свидетелями или публичными грамотами»}\footnote{Costums de Tortosa. 4.10.1: «En la ciutat de Tortosa prova hom tots feyts e tots contrayts, per tres coses: per confessió de son adversari, o per testimonis, o per cartes publiques». P.~202.}. Далее составители поясняли, что к письменным актам, способным выполнять эту функцию, относились публичные грамоты, завещания, кодициллы и иные последние волеизъявления, в которых участвовали два или три свидетеля и писец: «Письменными актами, то есть публичными [грамотами], либо завещаниями, кодициллами или иными последними волеизъявлениями, при которых присутствовали два или три свидетеля или более и писец»\footnote{Ibid. 4.10.1: «Per cartes: ço és a saber, publiques, o per testaments o codicils o altres derreres volentats en què aja dos o tres testimonis o plus e l’escrivà». P.~202.}.

Построение этой нормы показывает, что составители различали публичную грамоту и завещательные распоряжения, действительность которых требовала специального удостоверения свидетелями и писцом. Публичная грамота названа без дополнительного разъяснения её состава: по-видимому, сама принадлежность документа к категории публичных актов предполагала известный читателю порядок его составления. Завещание, кодицилл и иная последняя воля, напротив, нуждались в прямом указании на участников удостоверения. Из этого не следует, что публичная грамота вообще не была связана со свидетелями, однако её правовое значение выводилось прежде всего из публичной формы, тогда как для завещательного распоряжения источник отдельно фиксировал присутствие свидетелей и писца.

Роль писца в создании публичного акта не сводилась к технической записи слов сторон. Согласно статье 4.9.4, должник мог принять на себя обязательство посредством публичной грамоты даже в отсутствие кредитора, поскольку писец, перед которым составлялся документ, занимал его место: \emph{«Общий обычай состоит в том, что всякий человек может посредством публичной грамоты обязаться перед другим как должник, независимо от того, присутствует кредитор или отсутствует, поскольку писец, перед которым составляется грамота, всегда занимает место кредитора»}\footnote{Ibid. 4.9.4: «Ús e costuma general és que tot hom se pot obligar ab carta pública deutor a altre, jassia ço que·l creedor sia present o absent, per ço com l’escrivà en qui poder se fa la carta tota via té loc del creedor». P.~200.}. То же правило распространялось на продажи, дарения и иные договоры, кроме соглашений, которые по своему характеру требовали личного присутствия сторон.

Эта норма особенно важна для понимания публичного удостоверения. Писец не только фиксировал уже совершённое перед ним действие, но и обеспечивал возможность его действительного совершения при отсутствии одной из сторон. Составление грамоты «перед писцом» или \emph{«в его }\emph{власти» }(кат. --- \emph{en}\emph{ }\emph{poder}\emph{ de }\emph{l’escrivà}) создавало правовую ситуацию, в которой присутствие должностного лица восполняло отсутствие кредитора. Следовательно, значение публичной грамоты основывалось не только на письменной форме, но и на особом положении писца, способного представлять отсутствующую сторону при составлении акта.

Именно поэтому публичное удостоверение следует отличать от простого письменного изложения частного соглашения. В первом случае текст создавался при участии лица, за которым местный правовой порядок признавал полномочия составлять грамоты и обеспечивать их действие. Писец соединял волеизъявление стороны с установленной формой акта и тем самым позволял впоследствии ссылаться не только на содержание документа, но и на обстоятельства его официального составления.

Особые требования предъявлялись к завещаниям и иным последним волеизъявлениям. В «Обычаях Тортосы» действительность завещательной грамоты связывалась с записью, участием городского писца и формальным подтверждением со стороны свидетелей. Статья 6.4.26 приводила образец свидетельской подписи: \emph{«Я, такой-то, по просьбе такого-то подписываюсь как свидетель и ставлю свой знак»}\footnote{Ibid. 6.4.26: «Ego talis rogatus a tali me subscribo pro teste et signum meum appono». P.~314.}. В данном случае условия признания документа были выражены через совокупность внешних признаков: наличие записи, подписи или знака завещателя, подписей свидетелей и участие писца. Каждый из этих элементов связывал текст с конкретными лицами, принимавшими участие в его составлении и способными подтвердить надлежащее совершение распоряжения.

При этом значение писца состояло не только в его грамотности. Если завещатель не мог самостоятельно подписать документ, обращение к городскому писцу включало распоряжение в признанную форму правового письма. Поставленные свидетелями знаки также не являлись простым украшением текста: они обозначали личное участие в удостоверении и позволяли отличить надлежащим образом составленный акт от неподтверждённой записи. Таким образом, публичное или нотариальное удостоверение складывалось из взаимосвязанных признаков --- личности составителя, принятого формуляра, подписей и знаков участников.

Сходная модель получила дальнейшее развитие в привилегии «Знатные мужи Барселоны признали…» (Recognoverunt proceres). Согласно статье XXV, нотариус мог составить завещание, находясь наедине с завещателем, однако после записи текста на бумаге он должен был вызвать свидетелей и объявить перед ними, что составил завещание данного лица. Такое завещание признавалось действительным, как если бы свидетели сами слышали волеизъявление завещателя: \emph{«Нотариус может составить завещание, находясь наедине с завещателем, а после его составления или записи на бумаге должен призвать свидетелей, перед которыми объявит, что составил завещание этого завещателя; и оно имеет такую силу, как если бы свидетели сами слышали завещание»}\footnote{Recognoverunt proceres. [XXV]: «Item, quod notarius potest facere testamentum ipso solo existente cum testatore, et quod, ipso facto sive notato in papiro, vocet testes coram quibus dicat se fecisse testamentum ipsius testatoris, et quod valet ac si audivissent testamentum ipsi testes». P.~595.}.

Эта норма заметно изменяла распределение функций между нотариусом и свидетелями. Свидетели уже не обязательно присутствовали при первоначальном выражении последней воли. Они слышали последующее заявление нотариуса о том, что завещание было составлено. Именно признанное положение нотариуса позволяло приравнять такое опосредованное удостоверение к непосредственному восприятию слов завещателя. Достоверность акта связывалась не только с числом свидетелей, но и с ответственностью нотариуса за соответствие записи полученному волеизъявлению.

Вместе с тем нотариальная форма не была единственным способом придать последней воле правовое значение. Глава XLVIII той же привилегии признавала действительным завещательное распоряжение, совершённое письменно или устно в отсутствие нотариуса, если свидетели в течение шести месяцев являлись в Барселоне и приносили клятву в церкви Святого Юста\footnote{Ibid. [XLVIII]: «Item, est consuetudo quod si aliquis fecerit testamentum presentibus testibus, in terra vel in mare ubicumque sit, in scriptis vel sine scriptis, seu suam voluntatem etiam aliquo notario non presente, in ipsa voluntate verbotenus dicta vel scripta, quod valeat ipsa ultima voluntas sive testamentum, dum testes qui interfuerint ipsi ultime voluntati vel testamento infra sex menses ex quo fuerint in Barchinona iurent in ecclesia sancti Iusti, super altare sancti Felicis martiris». P.~599.}. Они должны были подтвердить, что видели и слышали, как завещание было записано или произнесено, и что последующая запись соответствует словам завещателя. Такое завещание называлось сакраментальным.

Сопоставление двух норм показывает, что в каталонской правовой культуре сосуществовало несколько способов удостоверения последней воли. В первом случае центральное положение занимал нотариус, чьё заявление восполняло отсутствие свидетелей при составлении текста. Во втором отсутствие нотариуса компенсировалось последующей клятвой свидетелей перед публично определёнными участниками и в установленном месте. И нотариальное завещание, и сакраментальное строились на включении частного волеизъявления в процедуру, которая позволяла публично подтвердить его происхождение и содержание.

Публичное удостоверение распространялось не только на акты, создававшие или передававшие имущественные права, но и на записи, возникшие в ходе судебной деятельности. В 51 статье привилегии «Знатные мужи Барселоны признали…» (Recognoverunt proceres) устанавливалось, что сообщению сайона о вызове стороны в суд и назначении дня заседания следовало доверять при условии, что соответствующая запись обнаруживалась в реестре или судебных актах: \emph{«Сообщению одного }\emph{сайона}\emph{ доверяют и следуют ему, если только запись обнаруживается в реестре или судебных актах»}\footnote{Ibid. [LI]: «Creditur et statur relationi unius sagionis, dummodo reperiatur scriptura in capibrevio vel in actis iudicis». P.~600.}.

Здесь официальная запись выполняла функцию проверки сообщения должностного лица. Доверие связывалось не только с положением сайона, но и с возможностью сопоставить его слова с записью, внесённой в реестр или судебные акты. Тем самым письменная фиксация становилась частью деятельности самого суда: она не просто сохраняла сведения о произведённом вызове, а позволяла признать сообщение о нём достаточным для дальнейшего движения дела.

Особенно отчётливо зависимость правового значения документа от места его составления выражена в «Кутюмах Перпиньяна». Согласно статье LXII, истец, требовавший выплаты денежной суммы на основании публичного документа, составленного в писчей канцелярии Перпиньяна, не был обязан подробно излагать основание требования: \emph{«Истец не обязан полностью излагать основание иска или требования, если требует уплаты причитающейся ему денежной суммы посредством публичного акта, составленного в писчей канцелярии Перпиньяна»}\footnote{Les Coutumes de Perpignan. [LXII]: «Nemo agens teneatur in tota causa exprimere causam seu peticionis, dum tamen eam de pecunia sibi solvenda faciat cum instrumento publico facto in scribania Perpinyani». P.~584--585.}.

Таким образом, происхождение документа из городской писчей канцелярии имело непосредственные процессуальные последствия. Публичный акт частично заменял развёрнутое изложение обстоятельств, из которых возникло требование: сторона могла опереться на документ как на уже оформленное основание иска. Значение имело не только то, что долг был записан, но и то, что запись была произведена именно в признанной городской канцелярии.

Та же связь между происхождением акта и порядком его признания проявлялась в статье XI. Против документов писчей канцелярии Перпиньяна не допускалось доказательство, кроме показаний свидетелей, названных в самих актах, либо показаний нотариуса и свидетелей: \emph{«Против актов писчей канцелярии Перпиньяна не допускается доказательство, кроме }\emph{показаний свидетелей этих актов либо нотариуса и свидетелей»}\footnote{Ibid. [XI]: «Contra instrumenta scribanie Perpinyani non admittitur probatio, nisi per testes eorum instrumentorum, vel nisi per tabellionem et testes». P.~570.}. Подробно порядок опровержения таких документов будет рассмотрен далее; здесь существенно, что принадлежность акта к документации городской канцелярии ограничивала круг средств, посредством которых его содержание могло быть поставлено под сомнение.

Материал Тортосы, Барселоны и Перпиньяна позволяет выделить несколько признаков публичного удостоверения. Во-первых, документ составлялся лицом, чьи полномочия признавались местным правовым порядком: писцом, нотариусом или служащим городской канцелярии. Во-вторых, акт включался в установленную процедуру, которая могла предусматривать запись в реестре, вызов свидетелей, подписи и знаки участников либо последующее публичное подтверждение. В-третьих, происхождение документа из канцелярии или нотариальной практики влекло определённые процессуальные последствия: позволяло действовать при отсутствии одной из сторон, сокращало необходимость повторного изложения основания требования или ограничивало способы опровержения записи.

При этом публичное удостоверение не превращало документ в безусловное и неоспоримое доказательство. Даже акт городской канцелярии сохранял связь с нотариусом и свидетелями, а завещательная запись могла требовать дополнительного подтверждения. Особое значение публичного акта состояло не в полном устранении других средств установления обстоятельств дела, а в том, что часть проверки переносилась с личности конкретного участника на порядок составления и хранения документа. Суду и сторонам уже не требовалось каждый раз заново устанавливать все обстоятельства совершённого действия: соблюдение признанной формы создавало исходное основание доверять записи до тех пор, пока она не была поставлена под сомнение в предусмотренном порядке.

Таким образом, к концу XIII~в. в каталонских городах формировался порядок, при котором правовое значение письменного акта всё теснее связывалось с деятельностью специальных участников и учреждений. Писец, нотариус, городская канцелярия и судебный реестр становились носителями устойчивости документа во времени. Это не устраняло свидетельское удостоверение, но изменяло его место: свидетели всё чаще подтверждали не только само правовое действие, но и надлежащее соблюдение процедуры составления акта. Публичная форма придавала документу особое положение в тяжбе, поскольку позволяла опереться на институционально закреплённую запись, а не только на память и добрую славу отдельных лиц.

\subsection{Предъявление и оспаривание письменного акта}

Письменный акт приобретал значение для разрешения тяжбы не вследствие одного лишь факта своего существования. Он должен был быть представлен суду, соотнесён с предметом требования и включён в состязание доводов сторон. В свою очередь, противная сторона могла попытаться опровергнуть содержание документа, противопоставить ему другой акт либо поставить под сомнение его достоверность. Поэтому правовое значение грамоты определялось не только условиями её составления, но и порядком её использования в ходе разбирательства.

В «Барселонских обычаях» процессуальная последовательность обращения с письменным актом наиболее ясно представлена в статье C3~(148). Одна из сторон могла предъявить в тяжбе завещание или надлежащим образом скреплённую грамоту, относившуюся к предмету спора. Противная сторона получала возможность опровергнуть предъявленный акт посредством свидетелей или других удостоверенных документов:\emph{ «Если кто-либо предъявит в тяжбе завещание или скреплённую грамоту относительно предмета спора, а другая сторона не сможет опровергнуть их ни свидетелями, ни надлежащим образом удостоверенными }\emph{письменными актами, судья должен постановить то, что представляется ему правильным, и сохранить за каждым его право»}\footnote{Usatges de Barcelona. Us. C3 (148): «Si quis testamentum uel cartam firmatam de aliqua contencione in placito ostenderit, de alia parte neque per testes neque per firmas scripturas conuincere potuerit, iudex dicat quod rectum ei uidetur, et seruet unicuique suum directum». P.~171.}.

Содержание нормы образуют три последовательных действия: предъявление документа одной стороной, попытка его опровержения другой и оценка представленных сведений судьёй. Глагол \emph{«предъявить»} (лат.~--- \emph{ostendere}) указывает, что документ должен был быть непосредственно введён в разбирательство, а не просто упомянут стороной как существующее основание права. При этом предмет грамоты должен был совпадать с предметом тяжбы: источник говорит именно об акте, составленном\emph{ «относительно спора»} (лат. --- \emph{de }\emph{aliqua}\emph{ }\emph{contencione}).

Опровержение могло строиться как на свидетельских показаниях, так и на других письменных актах. Следовательно, документ не только конкурировал со свидетельским словом, но и мог быть поставлен в один ряд со встречными грамотами. Норма не объясняла, по каким критериям судья должен был выбирать между противоречившими друг другу документами, однако обозначение встречных записей как надлежащим образом удостоверенных (лат. --- \emph{firmae}\emph{ }\emph{scripturae}) показывает, что их сила зависела от признанной формы составления. Простого письменного заявления для опровержения первоначальной грамоты было недостаточно.

Даже отсутствие убедительного опровержения не превращало предъявленный акт в автоматически обязательное основание решения. Судье предписывалось постановить то, что представлялось ему правильным, сохранив за каждым его право. Это не следует понимать как неограниченную свободу суда игнорировать документ. Скорее, статья показывала, что грамота входила в совокупность обстоятельств, подлежавших судебной оценке, и не заменяла самого решения. Письменный акт усиливал положение стороны, но его правовое последствие определялось только после рассмотрения возражений противника.

Сходная последовательность обнаруживается в актовой практике конца XII~в. В решении 1197~г. по спору о владениях Масана и Пуж судья Бернат де Бушо специально перечислил совершённые им действия: он выслушал доводы обеих сторон, рассмотрел и исследовал письменные акты, принял свидетелей и тщательно допросил их: «\emph{выслушав доводы и утверждения обеих сторон, рассмотрев и исследовав документы, приняв также свидетелей и тщательно расспросив их»}\footnote{Justícia i resolució…2024. Doc. 830: «Auditis racionibus et allegacionibus utriusque partis, visis instrumentis et cognitis, testibus etiam receptis et diligenter examinatis». P.~1023.}.

По итогам разбирательства спорные владения были присуждены монастырю Сан-Кугат в соответствии с документами, представленными монахами, и свидетельскими показаниями, подтверждавшими эти документы\emph{: «как содержится в актах, предъявленных монахами Сан-}\emph{Кугата}\emph{, и как подтверждают показания приведённых ими свидетелей»}\footnote{Ibid. Doc. 830: «Sicut continetur in instrumentis a monachis Sancti Cucufatis prolatis atque utpote dicta testium a predictis monachis prolata asserunt». P.~1024.}. В данном случае грамоты не действовали изолированно: суд сопоставлял их с доводами сторон и показаниями свидетелей.

Вместе с тем решение сохраняло за противной стороной возможность подтвердить часть своих требований в будущем. Дочери Арнау де Фонтайата могли получить права на приобретения Гильема Отона, если сумели бы доказать посредством документов или свидетелей своё происхождение от него и наличие соответствующего наследственного распоряжения. Местонахождение приобретённых владений также следовало установить\emph{ «посредством документов, свидетелей или описаний границ»}\footnote{Ibid. Doc. 830: «Per instrumenta vel testes seu per affrontaciones». P.~1024.}.

Этот казус показывает, что предъявление документа не было одномоментным действием, после которого суд только подтверждал записанное в нём право. Письменные акты исследовались, сопоставлялись с показаниями и использовались для установления разных обстоятельств: принадлежности имущества, родства, содержания наследственного распоряжения и местоположения владений. Один и тот же документ мог быть достаточным для подтверждения определённого элемента требования, но не устранять необходимости отдельно установить другие обстоятельства.

«Обычаи Льейды» более отчётливо регулировали возможность поставить письменный акт под сомнение. Рубрика «О доверии к письменным актам» устанавливала: \emph{«Любой письменный акт может быть объявлен подозрительным, даже если на нём нет соответствующей надписи}\emph{,}\emph{ и он не отменён или не стёрт»}\footnote{Costumbres de Lérida. [134]: «Suspectum potest quis dicere quodlibet instrumentum etiam sine inscripcione non abolitum vel deletum». P.~188.}. Далее следует формула \emph{«и этому надлежит придать доверие»}\footnote{Ibid. [134]: «Et est ei fides facienda». P.~188.}, синтаксическая связь которой с предшествующей частью нормы не вполне ясна.

Поэтому из этой статьи нельзя без дополнительных оснований заключать, что одного заявления о подозрительности было достаточно для признания документа недействительным. Более осторожно можно утверждать, что сторона имела право оспорить акт даже при отсутствии явных внешних признаков его отмены или повреждения, а суд не должен был отвергать такое заявление исключительно потому, что документ сохранял обычный внешний вид. Объявление акта подозрительным открывало его проверку.

В пользу такой интерпретации говорит и общий порядок оценки документов в «Обычаях Льейды». Показание одного свидетеля не принималось против письменного акта в качестве достаточной презумпции\footnote{Ibid. [129]: «Unius testis presumpcio contra instrumentum non admittitur». P.~187.}. Это правило защищало документ от слабого возражения, но не исключало его опровержения более полным подтверждением. Таким образом, право одновременно допускало заявление о подозрительности грамоты и предъявляло повышенные требования к средствам, которыми это сомнение должно было быть обосновано.

Кроме того, суд мог отказаться приводить в исполнение отдельное условие письменного акта, не отрицая действительности документа в целом. Статья 136 «Обычаев Льейды» предусматривала, что записанное в долговой грамоте обязательство выплачивать дополнительную сумму за продление срока не должно взыскиваться судом, если представляло собой ростовщический процент: \emph{«Такое условие судья не принуждает исполнять»}\footnote{Ibid. [136]: «Illum talionem non facit iudex solvi». P.~188.}.

Следовательно, оспаривание могло относиться не только к подлинности или достоверности документа, но и к допустимости закреплённого в нём требования. Наличие условия в письменном акте само по себе не обязывало суд обеспечить его исполнение. Судья отделял существование записи от правового значения её содержания и мог признать отдельное положение акта не подлежащим принудительному осуществлению.

Иной порядок оспаривания был предусмотрен для документов писчей канцелярии Перпиньяна. Как было показано выше, их особое положение определялось происхождением из признанного городского учреждения. Процессуальным следствием этой формы было ограничение средств, допускавшихся против документа: \emph{«Против актов писчей канцелярии Перпиньяна не допускается доказательство, кроме показаний свидетелей этих актов либо нотариуса и свидетелей»}\footnote{Les Coutumes de Perpignan. [XI]: «Contra instrumenta scribanie Perpinyani non admittitur probatio, nisi per testes eorum instrumentorum, vel nisi per tabellionem et testes». P.~570.}.

Норма не делала канцелярский акт совершенно неоспоримым. Однако поставить его содержание под сомнение могли не любые свидетели, сообщавшие известные им обстоятельства, а лица, непосредственно связанные с составлением документа: свидетели, названные в нём, либо нотариус вместе со свидетелями. Тем самым проверка возвращалась к моменту создания акта. Оспаривающая сторона должна была показать, что именно участники удостоверения способны подтвердить несоответствие записи совершённому действию или нарушение порядка составления документа.

Связь проверки документа с лицами, присутствовавшими при его составлении или знавшими его содержание, прослеживается и в более раннем каталонском материале, рассмотренном Дж.А.~Боуменом. На рубеже X--XI~вв. показания таких лиц могли не только подтверждать или опровергать предъявленную грамоту, но и служить основанием для восстановления утраченного документа: содержание акта воспроизводилось по памяти тех, кто видел его и был знаком с записанным в нём распоряжением\footnote{\emph{Bowman J.A.}~Shifting Landmarks: Property, Proof, and Dispute in Catalonia around the Year 1000. Ithaca; London: Cornell University Press, 2004. P.~152--156.}. Письменный акт и свидетельское слово в подобных случаях не выступали как взаимоисключающие средства подтверждения права, поскольку показания обеспечивали сохранение действия документа при утрате его материального носителя. Норма «Кутюмов Перпиньяна» регулировала иную ситуацию, однако также связывала проверку акта с участниками его удостоверения. Отличие заключалось в более строгом определении их круга: вместо потенциально более широкого сообщества лиц, хранивших память о документе, правовой памятник называл свидетелей, указанных в самом акте, и нотариуса.

Сопоставление «Барселонских обычаев», «Обычаев Льейды» и «Кутюмов Перпиньяна» позволяет увидеть несколько способов регулирования спора вокруг письменного акта. В первом памятнике предъявленная грамота могла быть опровергнута свидетелями или другими удостоверенными документами, после чего окончательную оценку давал судья. В Льейде любой акт мог быть объявлен подозрительным, но отдельные нормы одновременно защищали письменную запись от недостаточно убедительных возражений и позволяли суду отказаться от исполнения неправомерного условия. В Перпиньяне публичное происхождение документа существенно сужало круг допустимых средств его опровержения.

Даже публичный акт городской канцелярии мог быть поставлен под сомнение, тогда как более ранняя скреплённая грамота уже могла обеспечить стороне сильную позицию в тяжбе. Изменялся прежде всего порядок проверки: чем теснее документ был связан с публично признанной процедурой составления, тем более определёнными становились средства и лица, посредством которых допускалось его оспаривание.

Таким образом, письменный акт не устранял состязательного характера разбирательства. Сторона должна была предъявить его и показать связь с предметом требования; противник мог противопоставить свидетелей, другой акт, заявление о подозрительности либо возражение против содержания записанного условия. Возрастание значения письменной формы выражалось поэтому в постепенной регламентации порядка предъявления, проверки и опровержения грамоты.

\subsection{Утрата, восстановление и доказательная сила документа}

Предъявление и проверка письменного акта разрешали вопрос о его значении в конкретной тяжбе, однако не исчерпывали проблемы сохранения подтверждённого им права. Материальный носитель мог быть утрачен, повреждён или уничтожен, а завершённый конфликт --- возобновиться после смены поколений либо участников отношений. Поэтому письменная практика должна была обеспечивать не только удостоверение права в момент составления акта, но и его сохранение во времени. В актовом материале обнаруживаются два взаимосвязанных способа решения этой задачи. Поэтому необходимо различать письменное закрепление достигнутого решения, восстановление утраченного акта и проверку документа, подлинность либо действительность которого ставилась под сомнение: каждая из этих ситуаций предполагала собственное соотношение письменной записи, свидетельства и судебной оценки.

Наиболее ранний и развёрнутый пример восстановления утраченных актов связан с гибелью монастырских документов в Эшаладе при разливе реки Тет. Представленные монастырём свидетели приносили клятву относительно грамот о покупке, дарении и обмене, утраченных во время наводнения\footnote{Justícia i resolució... 2018. Doc. 25: «…scripturae emptionis, donationis et commutationis, quae perditae fuerunt in Exalata ad delavatione fluvii Tete». P.~73.}. Они утверждали, что лично видели эти документы, слышали, как их неоднократно читали, и присутствовали при их публичном оглашении: \emph{«Мы видели и слышали, как эти грамоты читали и перечитывали»}\footnote{Ibid. Doc. 25. P.~73.}.

Как было показано выше, свидетельское показание в данном случае восполняло отсутствие первоначального документа. Однако для рассматриваемого здесь сюжета существенно другое: восстановление права не ограничивалось устным сообщением о существовании утраченных грамот. Показания были включены в новый письменный акт --- условия клятвы свидетелей, составленные по распоряжению судей. В результате сведения, прежде содержавшиеся в нескольких документах о приобретении и передаче имущества, получили новую форму письменного удостоверения. Устная память свидетелей не просто заменяла утраченный текст в ходе одного заседания, а становилась основанием для создания записи, способной сохранять восстановленные сведения и после завершения разбирательства.

Различие между умышленным устранением документа и его случайной утратой было проведено уже во «Вестготской правде». Статья VII.5.2 устанавливала ответственность за составление и предъявление ложной записи, сокрытие, похищение или повреждение подлинной грамоты, но одновременно допускала восстановление документа, погибшего вследствие небрежности, пожара, наводнения или иного несчастного случая\footnote{LV. VII.5.2. P.~304--305.}.

Состав нового документа не был буквальной копией погибших грамот. Свидетели воспроизводили известные им сведения о совершённых покупках, дарениях и обменах, об участниках этих действий и соответствующем имуществе. Рассматривая каталонские случаи восстановления утраченных записей, Дж.А.~Боумен связал эту процедуру с взаимодействием письменного акта, свидетельского слова и памяти лиц, видевших документ или присутствовавших при его чтении\footnote{\emph{Bowman J.A. }Shifting Landmarks: Property, Proof, and Dispute in Catalonia around the Year 1000. Ithaca; London: Cornell University Press, 2004. P.~152--156.}.

Такой порядок позволял разграничить утрату материального носителя и прекращение подтверждённого им права. Гибель грамоты сама по себе не означала, что приобретение или дарение переставали существовать. Однако стороне требовалось восстановить сведения, которые позволяли установить содержание утраченного акта и связь с ним заявленного требования. Судебное свидетельство и его письменная фиксация обеспечивали продолжение действия права при отсутствии первоначального документа.

«Барселонские обычаи» и иные локальные своды обычного права не предусматривали специального наказания за подлог документов. При этом в арагонских сводах обычного права, например, в «Фуэро Теруэля» санкции за это правонарушение были обозначены эксплицитно\footnote{\emph{Iglesias }\emph{Rábade L.} Estudio comparado del régimen jurídico del delito de falsedad documental en el Derecho hispánico e inglés en el Medievo // Estudios de Deusto. 2016. Vol.~64, №~2. P.~76. DOI: 10.18543/ed-64(2)-2016pp67-100.}. Ряд исследователей предположил, что в данном случае в Каталонии даже в XII--XIII~вв. продолжали пользоваться нормами «Вестготской правды». Как нам кажется, неразработанность норм, связанных с ответственностью за подлог документов, могла быть связана с высоким уровнем документальной культуры и хорошей сохранностью оригинальных грамот\footnote{\emph{Idem.} Op. cit. P.~76; \emph{Alejandre García J.A.~Estudio histórico del delito de falsedad documental // Anuario de Historia del Derecho Español. 1972. T.~42. P.~117--188. Здесь: P.~153--156; }\emph{Филиппов И.С.} Рукописи не горят: проблема реконструкции погибших архивов и возможности расширения базы источников по истории Средневековья // Как мы пишем историю? / отв. ред. Г. Гаррета, Г. Дюфо, \emph{Л.А.~Пименова}. М.: РОССПЭН, 2013. С.~213--214; \emph{Сафронова Д.П. }Фальшивые грамоты в леонской правовой практике X--XI веков // История: переводить, понимать, оценивать (к юбилею \emph{М.А.~Юсима}) / отв. ред. \emph{П.Ю.~Уваров}. М.: Институт всеобщей истории РАН, 2021. С.~123--124.} в отличие от других регионов, где зачастую оригинальные документы не сохранялись, и в распоряжении современников были в основном сборники копий документов --- картулярии.

Вместе с тем использование свидетельской памяти предполагало специальную проверку самого факта утраты. В подобных актах лицо, в пользу которого восстанавливался документ, должно было подтвердить, что не скрывает первоначальную грамоту и не заявляет о её исчезновении обманным путём. Иначе существовала бы опасность одновременного обращения и с первоначальным актом, и с записью, созданной для его замены. Поэтому восстановление документа соединяло подтверждение его содержания с удостоверением обстоятельств утраты.

Восстановление случайно утраченного акта следует отличать от изготовления, изменения или использования ложной грамоты. В первом случае процедура была направлена на сохранение ранее возникшего права при исчезновении его письменного подтверждения; во втором речь шла о сознательном создании или использовании записи, способной исказить основания притязаний одной из сторон. Различались и последствия: восстановление завершалось составлением нового документа, тогда как выявленный подлог мог повлечь утрату актом силы и наказание виновного.

Необходимость сохранять право при возможной гибели грамот приводила и к созданию вспомогательных записей, содержавших сведения о нескольких документах. Такие перечни не являлись полными копиями каждого акта, но фиксировали их существование, вид и относящееся к ним имущество. Их назначение состояло в том, чтобы сохранить письменную память о совокупности оснований права и облегчить последующее восстановление сведений при утрате отдельных грамот. В отличие от записей условий клятвы такие документы могли создаваться заранее, ещё до возникновения тяжбы, как средство защиты от возможной потери документов.

Иная задача возникала после разрешения конфликта. Судебное решение, признание требования или достигнутое примирение должны были сохранять силу не только для непосредственных участников, но и для их наследников и преемников. Эту функцию выполняли грамоты об окончательном отказе от притязаний, признании права другой стороны и прекращении спора. В актовом материале для них употреблялись обозначения «окончательный отказ» (лат. --- \emph{diffinitio}), «уступка или освобождение от притязаний» (лат. --- \emph{evacuatio}), «признание» (лат. --- \emph{recognitio}) и «примирение» (лат. --- \emph{pacificatio}).

Акты, оформлявшие отказ от притязаний и прекращение тяжбы, были рассмотрены в главе 2 с точки зрения их формулярного и речевого устройства. Здесь существенно иное: составленный по итогам разбирательства документ закреплял признанное сторонами распределение прав и мог быть предъявлен позднее как подтверждение того, что прежние требования были прекращены. Однако для большинства актов решающим становилось включение отказа в устойчивый формуляр. Письменная формула позволяла представить действие как завершённое, окончательное и распространяющееся на будущее. Именно благодаря этому грамота могла впоследствии быть предъявлена против новой попытки заявить уже прекращённое требование.

Следует различать функцию подобных актов и доказательную силу обычной грамоты о продаже или дарении. Последняя подтверждала действие, совершённое до возникновения конфликта, и могла служить основанием требования. Грамота об окончательном отказе, напротив, создавалась вследствие уже возникшего спора и закрепляла изменение правового положения сторон. Она не только доказывала, что одна из сторон признала право другой, но и устанавливала невозможность вернуться к прежней позиции без нарушения принятого обязательства.

Каталонский актовый материал показывает, что предъявление ложной грамоты могло влечь предусмотренное законом взыскание. В документе 980~г. епископ Жироны Миро упоминал жителей Палау-де-Санта-Эулалия, которые представили в суде графа Госфреда I \emph{«совершенно }\emph{ложную грамоту»}. При этом епископ освобождал их от заслуженной ими \emph{«кары нашего закона»}\footnote{Justícia i resolució… 2024. Doc. 89: «…qui scripturam falsissimam in eius iuditio protulerunt <…> nostreque legis censura qua merebamini eis indulgemus». P.~195.}. Составители этой грамоты не называли конкретной нормы о подлоге, поэтому связывать санкцию непосредственно с VII.5.2 «Вестготской правды» можно лишь предположительно. Однако сам казус показывает, что ложный документ не просто утрачивал силу в тяжбе: его предъявление рассматривалось как деяние, заслуживавшее самостоятельного наказания.

Разрыв между разработанностью нормативных предписаний и способами проверки документов в конкретных тяжбах прослеживается и в леонском материале X--XI~вв., исследованном Д.П.~Сафроновой\footnote{\emph{Сафронова Д.П.} Фальшивые грамоты в леонской правовой практике X--XI веков // История: переводить, понимать, оценивать (к юбилею \emph{М.А.~Юсима}) / отв. ред. \emph{П.Ю.~Уваров}. М.: Институт всеобщей истории РАН, 2021. С.~139--140.}. Наряду с дипломатическими и историческими подделками она рассматривает и иные махинации с грамотами, включая составление акта под принуждением и использование формально оформленного документа для представления недобровольной передачи имущества как законной сделки. Хотя \emph{Liber}\emph{ }\emph{Iudiciorum} предусматривал суровые санкции и специальные способы выявления подлога, в рассмотренных исследовательницей конфликтах сомнительные грамоты чаще опровергались клятвой и свидетельскими показаниями. Выявленный ложный документ могли признать недействительным и разорвать, однако это не всегда окончательно прекращало имущественные притязания представившей его стороны.

Следовательно, значение документа состояло не только в способности подтвердить обстоятельства уже начавшейся тяжбы. Письменный акт связывал различные моменты во времени: сохранял память о ранее совершённом действии, позволял восстановить утраченные сведения и закреплял последствия разрешённого конфликта для наследников и преемников. При этом устойчивость записи продолжала поддерживаться свидетельской памятью, публичным удостоверением, санкциями и дополнительными гарантиями исполнения. Письменная форма не устраняла эти элементы, но соединяла их в акте, рассчитанном на продолжительное действие.

\section*{Выводы к главе 3}
\addcontentsline{toc}{section}{Выводы к главе 3}

Проведённый анализ показывает, что судебное доказательство в Каталонии XII--XIII~вв. не было организовано в виде единой и неизменной иерархии. «Барселонские обычаи», городские привилегии и локальные сборники закрепляли различные сочетания ордалии, судебного поединка, клятвы, свидетельских показаний, репутационных критериев и письменных актов. Их применение зависело от характера дела, статуса участников и порядка, принятого в конкретном городе. Изменения затрагивали не только распространённость отдельных способов доказательства, но прежде всего условия их допустимости, круг лиц, участвовавших в их совершении, и последствия полученного результата.

Ордалия и судебный поединок в «Барселонских обычаях» представляли собой установленные и подробно регулируемые способы разрешения спора. Выбор испытания зависел от социального положения участника, его пола, имущественных и физических возможностей. Рыцари должны были участвовать в поединке с равными себе, женщины привлекали представителей в соответствии со статусом семьи, тогда как крестьянки лично проходили испытание кипятком. Размер залога, порядок возмещения ущерба и имущественные последствия победы или поражения также определялись заранее. Таким образом, судебное испытание соединяло установление правоты с признанной социальной иерархией сторон.

Актовый материал XI--XII~вв. подтверждает особое положение судебного поединка среди видов «Божьего суда». Он встречался заметно чаще односторонних испытаний водой или железом и сохранял значение при разрешении различных конфликтов. В XIII~в. область применения ордалий стала сокращаться, однако этот процесс протекал неодинаково. «Обычаи Льейды» исключали подобные испытания, в Тарреге они допускались по взаимному согласию сторон, а в Тортосе судебный поединок и иные ордалии не признавались местным порядком. В Барселоне поединок сохранялся применительно к отдельным тяжким деяниям. Поэтому речь шла не об одновременном исчезновении «Божьего суда», а о различавшемся в отдельных городах ограничении его применения.

Клятва занимала промежуточное положение между самостоятельным способом доказательства и обязательным элементом судебной процедуры. В «Барселонских обычаях» она могла предшествовать ордалии или судебному поединку, дополнять их либо заменять в пределах установленной суммы. В «Обычаях Тортосы» клятва сохраняла способность непосредственно влиять на исход дела, но одновременно включалась в последовательность процессуальных действий: сопровождала заявление стороны, предшествовала даче показаний и подтверждала отказ от дальнейшего доказательства. Отказ принести назначенную клятву мог быть приравнен к признанию позиции противника.

Это изменение отразилось и в терминологии. Глагол «клясться» (старокат. --- jurar) обозначал произнесение словесной формулы, тогда как конструкция «принести клятву» (старокат. --- fer sagrament) охватывала действие в целом. Другие сочетания с существительным sagrament --- «дать», «принять» или «возвратить клятву» --- передавали распределение процессуальной инициативы между сторонами. Клятва, следовательно, понималась не только как высказывание, но и как действие, правовые последствия которого зависели от соблюдения установленного порядка.

Особые правила для дел с участием иноверцев не создавали отдельного судебного порядка, но изменяли ритуальную форму клятвы, место её принесения и состав присутствующих лиц. В отношении иудеев «Обычаи Тортосы» предусматривали принесение клятвы на Моисеевом законе, использование проклятий и в отдельных случаях совершение соответствующих действий в синагоге. Такие установления позволяли учитывать конфессиональные различия, сохраняя общую процессуальную функцию клятвы. Сокращение подобных предписаний во второй редакции памятника показывает, что состав судебного ритуала мог пересматриваться без отказа от самой клятвы как правового действия.

Репутация определяла предварительное доверие к слову участника процесса. В «Барселонских обычаях» эта зависимость проявлялась прежде всего через ограничения, налагаемые на клятвопреступника: утративший доверие человек не мог впоследствии свидетельствовать или подкреплять свои слова клятвой. В «Обычаях Тортосы» добрая и дурная слава получили более подробное правовое выражение. От репутации зависели возможность предъявить обвинение, допустимость свидетельства, применение пытки и участие лица в розыскном разбирательстве.

Добрая слава не доказывала отдельное обстоятельство дела, но определяла пригодность человека к совершению процессуально значимого действия. Напротив, дурная слава ограничивала доверие к его слову и могла допускать применение принуждения. При этом репутация оценивалась применительно к моменту свидетельства или составления документа: позднейшее бесчестие не уничтожало силы показания, данного тогда, когда лицо ещё пользовалось доверием. Репутационный критерий был связан не только с нравственной оценкой поведения, но также с социальным положением, занятием, зависимостью и отношением к сторонам.

В «Обычаях Тортосы» дурная слава могла относиться и к городской общине. Нераскрытые преступления наносили ущерб не только отдельным пострадавшим, но и представлению о городе как о пространстве, в котором осуществляется правосудие. В таком случае расследование было направлено также на восстановление доверия к городскому порядку. Индивидуальная и коллективная репутация, таким образом, объединялись через представление о праве как средстве сохранения общественно признанного порядка.

Правила допуска свидетелей конкретизировали основания доверия к их словам. В различных памятниках устанавливались необходимое число свидетелей, возрастные ограничения, требования к репутации и отсутствие вражды или зависимости от одной из сторон. Предпочтение могло отдаваться лицам, принадлежавшим к той же местности и способным сообщить непосредственно воспринятые обстоятельства. Показания с чужих слов ограничивались, свидетелей предписывалось расспрашивать по отдельности, а в более подробных процедурах их слова записывались и оглашались обвиняемому.

Усложнение порядка допроса не устраняло персонального характера свидетельства. Правовое значение имело не только содержание сообщения, но и положение говорящего, характер его знания и связь с участниками разбирательства. Даже требование двух или трёх свидетелей не сводилось к количественному правилу: каждый из них должен был обладать процессуальной пригодностью. Вместе с тем свидетельство всё более определённо отделялось от обвинения и судебного решения. Разграничение функций обвинителя, свидетеля и судьи, правила вызова и порядок записи показаний формировали последовательность действий, ограничивавшую возможность произвольного использования свидетельского слова.

Письменные акты различались прежде всего по времени и цели составления. Грамоты о продаже, дарении, долге, передаче феода, брачном имуществе и последней воле возникали до конфликта и позднее могли подтверждать заявленное требование. Перечни жалоб, записи свидетельских показаний, судебные решения и обязательства сторон составлялись после начала разбирательства и закрепляли отдельные его действия. Документы об отказе от притязаний, признании права или примирении фиксировали достигнутый результат. Один и тот же акт мог менять свою функцию: грамота, первоначально удостоверявшая передачу имущества, становилась основанием требования, а документ, составленный после разрешения тяжбы, использовался при попытке её возобновления.

Распространение письменных актов не устранило значения свидетельского слова. Документ и показания могли служить самостоятельными основаниями требования, однако во многих случаях свидетели участвовали в составлении самого акта. Указание их имён позволяло вызвать этих лиц после возникновения спора и потребовать подтверждения совершённого при них действия. Вместе с тем документ мог удостоверять первоначальное обязательство, тогда как свидетели подтверждали его последующее исполнение. Письменная запись и устное сообщение относились в таком случае к разным этапам одного правового отношения.

Особенно ясно их взаимосвязь проявлялась при утрате грамоты. Содержание погибшего документа восстанавливалось по показаниям лиц, видевших его, присутствовавших при чтении или знавших удостоверительную часть. Одновременно лицо, заинтересованное в восстановлении акта, должно было подтвердить сам факт утраты. Новая запись не являлась буквальной копией прежнего текста, но воспроизводила признанные сведения о совершённом действии, его участниках, имуществе и порядке удостоверения. Гибель материального носителя поэтому не прекращала подтверждённого им права, если содержание акта могло быть установлено предусмотренным способом.

В городских памятниках XIII~в. особое значение приобретало публичное удостоверение. Правовая сила документа всё теснее связывалась с участием писца или нотариуса, происхождением записи из городской канцелярии и возможностью обнаружить её в реестре или судебных актах. В Тортосе писец не только записывал волеизъявление сторон, но в определённых случаях восполнял отсутствие одного из участников. В Барселоне признанное положение нотариуса позволяло составить завещание без непосредственного присутствия свидетелей при первоначальном распоряжении. В Перпиньяне публичный акт городской канцелярии освобождал истца от подробного изложения основания денежного требования.

Публичное происхождение документа не делало его безусловно неоспоримым. Грамота могла быть сопоставлена со свидетельскими показаниями и другими актами, объявлена подозрительной либо поставлена под сомнение лицами, участвовавшими в её удостоверении. Особое положение публичного документа состояло в том, что порядок его оспаривания был ограничен заранее. Проверка относилась уже не только к памяти сторон или свидетелей, но также к полномочиям составителя, соблюдению формуляра, наличию подписей и знаков, происхождению записи и условиям её хранения.

Следует также различать оспаривание документа, признание его недействительным и подлог. Акт мог правильно передавать совершённое действие, но утратить значение вследствие последующих обстоятельств. Он мог быть составлен при нарушении условий свободного волеизъявления либо содержать сведения, не соответствовавшие действительному характеру передачи имущества. Наконец, ложная грамота могла быть изготовлена, изменена или сознательно предъявлена в суде. В последнем случае речь шла не только об отказе признать документ основанием требования, но и об ответственности лица, использовавшего подлог.

Таким образом, изменения в области судебного доказательства состояли не в последовательном вытеснении ритуальных и устных форм письменными. Ордалия ограничивалась определёнными категориями дел и участников; клятва сохраняла сакральный характер, но включалась в более подробно установленную процедуру; репутация определяла допустимость слова; свидетельство подчинялось правилам отбора, допроса и записи; письменный акт получал особое положение благодаря порядку его удостоверения и хранения. Эти средства продолжали взаимодействовать, но их функции и соотношение становились более определёнными.

Рассмотрение судебного доказательства дополняет результаты предшествующих глав. Текстуальные связи между правовыми памятниками сопровождались изменением сферы применения заимствованных или унаследованных норм, а выбор языковой формы был связан с порядком совершения предусмотренного действия. Совокупность этих наблюдений позволяет рассматривать каталонскую правовую культуру XII--XIII~вв. через взаимосвязь состава правовых текстов, способов выражения нормы и организации судебного доказательства.
